% Paper v1 (MVP1): Q1 calibration, Q2 wallet concentration, Q3 descriptive. Run 142082.
% Every measured number is a macro from numbers.tex (generated, do not edit).
%
% Source key for the "% src:" comments (this key is the Sources table of this file):
%   S-PAPER-1 proposal_md/polymarket_proposal_v3.md
%   S-PAPER-2 tips/polymarket_plan_detailed.md
%   S-PAPER-3 analysis/prereg_mvp1.md
%   S-PAPER-5 runs/142082_mvp1_full_v1r3 (reports, verdicts.json) via numbers.tex
%   S-PAPER-6 analysis/review_mvp1_mvp2.md
%   S-PAPER-7 tips/polymarket_class_questions_tips.html
%   S-PAPER-8 PLAN.md, RQ.md, DECISIONS.md, STATUS.md
%   S-PAPER-9 analysis/prereg_robustness.md, runs/142123_robust_full_v4 (reports r1, r2) via numbers.tex
%   S-PAPER-10 reports/lecture_numbers.json, runs/145892_ccdf_v6, runs/145893_walletdist_v7 via numbers.tex; analysis/review_project_end_analyst.md
%   S-PAPER-11 analysis/related_work_refs.md (R1-R5)
%   S-PAPER-12 analysis/prereg_review_fixes.md (commit aef0808), runs/150310_r3_review_fixes_v6 via numbers.tex
\documentclass[11pt]{article}

\usepackage[margin=0.9in]{geometry}
\usepackage[T1]{fontenc}
\usepackage{lmodern}
\usepackage{microtype}
\usepackage{amsmath}
\usepackage{booktabs}
\usepackage{tabularx}
\usepackage[numbers]{natbib}
\usepackage{figkit}
\usepackage{subcaption}
\usepackage[hidelinks]{hyperref}
\usepackage[capitalise,noabbrev]{cleveref}

\captionsetup{font=small,labelfont=bf}
\setlength{\emergencystretch}{2em}
\setlength{\intextsep}{6pt plus 2pt minus 2pt}
\setlength{\textfloatsep}{8pt plus 2pt minus 2pt}
\setlength{\floatsep}{6pt plus 2pt minus 2pt}
\setlength{\abovecaptionskip}{4pt}

% Visible placeholder for text or numbers the sources do not give.
\DeclareRobustCommand{\todo}[1]{\textcolor{red}{\textbf{[TODO: #1]}}}

% Generated by code/make_paper_numbers.py from runs/142082_mvp1_full_v1r3. Do not edit by hand.
% run directory name
\newcommand{\RunId}{142082\_mvp1\_full\_v1r3}
% manifest_q1_calibration.json: code_version
\newcommand{\CodeVersion}{v3}
% manifest_q1_calibration.json: code_sha (hash of the script q1_calibration.py alone, first 12 hex, not of the code set)
\newcommand{\CodeSha}{8efcddf9ef94}
% code.sha256: sha256 of the file that lists the hash of every script (first 12 hex)
\newcommand{\CodeSetSha}{950ec496595b}
% path of the per-script hash list
\newcommand{\CodeShaFile}{runs/142082\_mvp1\_full\_v1r3/code.sha256}
% manifest_q1_calibration.json: env.slurm_job_id
\newcommand{\SlurmJobId}{142082}
% verdicts.json: comparisons.Q1.verdict
\newcommand{\VerdictQOne}{Prices are not calibrated as a whole at this n.}
% verdicts.json: alpha
\newcommand{\Alpha}{\ensuremath{0.05}}
% tables_summary.json: n_markets_in_scope
\newcommand{\NMarketsInScope}{\ensuremath{688{,}131}}
% tables_summary.json: n_kept
\newcommand{\NMarketsKept}{\ensuremath{47{,}948}}
% tables_summary.json: n_kept / n_markets_in_scope
\newcommand{\PctMarketsKept}{\ensuremath{7.0}\%}
% tables_summary.json: n_kept_no_min_filter
\newcommand{\NMarketsKeptNoMin}{\ensuremath{74{,}474}}
% tables_summary.json: drop_reasons.stale_price
\newcommand{\NStalePrice}{\ensuremath{80{,}230}}
% tables_summary.json: drop_reasons.stale_price / n_markets_in_scope
\newcommand{\PctStalePrice}{\ensuremath{11.7}\%}
% tables_summary.json: drop_reasons.below_min_trades
\newcommand{\NBelowMinTrades}{\ensuremath{26{,}526}}
% tables_summary.json: drop_reasons.below_min_trades / n_markets_in_scope
\newcommand{\PctBelowMinTrades}{\ensuremath{3.9}\%}
% tables_summary.json: drop_reasons.updown
\newcommand{\NUpdown}{\ensuremath{255{,}420}}
% tables_summary.json: drop_reasons.updown / n_markets_in_scope
\newcommand{\PctUpdown}{\ensuremath{37.1}\%}
% tables_summary.json: drop_reasons.kept
\newcommand{\NKept}{\ensuremath{47{,}948}}
% tables_summary.json: drop_reasons.kept / n_markets_in_scope
\newcommand{\PctKept}{\ensuremath{7.0}\%}
% tables_summary.json: drop_reasons.no_outcome
\newcommand{\NNoOutcome}{\ensuremath{62{,}718}}
% tables_summary.json: drop_reasons.no_outcome / n_markets_in_scope
\newcommand{\PctNoOutcome}{\ensuremath{9.1}\%}
% tables_summary.json: drop_reasons.no_price
\newcommand{\NNoPrice}{\ensuremath{215{,}289}}
% tables_summary.json: drop_reasons.no_price / n_markets_in_scope
\newcommand{\PctNoPrice}{\ensuremath{31.3}\%}
% tables_summary.json: n_trades_in_kept_markets
\newcommand{\NTradesKeptMarkets}{\ensuremath{63{,}023{,}987}}
% tables_summary.json: n_wallet_role_rows
\newcommand{\NWalletRoleRows}{\ensuremath{2{,}117{,}412}}
% tables_summary.json: trade_size_sample_rows
\newcommand{\NTradeSizeSample}{\ensuremath{1{,}000{,}000}}
% tables_summary.json: n_kept_open_ge_7d
\newcommand{\NKeptOpenSevenDays}{\ensuremath{21{,}816}}
% tables_summary.json: n_kept_with_p_7d
\newcommand{\NKeptWithPriceSevenDays}{\ensuremath{14{,}169}}
% q1_results.json: main_with_min_filter.n_markets
\newcommand{\QOneN}{\ensuremath{47{,}948}}
% q1_results.json: main_with_min_filter.chi2_gof.chi2
\newcommand{\QOneChiSq}{\ensuremath{68.0}}
% q1_results.json: main_with_min_filter.chi2_gof.df
\newcommand{\QOneDf}{\ensuremath{10}}
% q1_results.json: main_with_min_filter.chi2_gof.p
\newcommand{\QOneP}{\ensuremath{1.1\times10^{-10}}}
% q1_results.json: main_with_min_filter.chi2_gof.effect_w
\newcommand{\QOneW}{\ensuremath{0.038}}
% q1_results.json: main_with_min_filter.brier
\newcommand{\QOneBrier}{\ensuremath{0.107}}
% q1_results.json: main_with_min_filter.log_loss
\newcommand{\QOneLogLoss}{\ensuremath{0.323}}
% q1_results.json: main_with_min_filter.mean_entropy_bits
\newcommand{\QOneEntropy}{\ensuremath{0.475}}
% q1_results.json: main_with_min_filter.class_balance_yes
\newcommand{\QOneYesRate}{\ensuremath{36.6}\%}
% q1_results.json: main_with_min_filter.confusion_at_0.5.accuracy
\newcommand{\QOneAccuracy}{\ensuremath{83.4}\%}
% q1_results.json: main_with_min_filter.confusion_at_0.5.precision
\newcommand{\QOnePrecision}{\ensuremath{77.7}\%}
% q1_results.json: main_with_min_filter.confusion_at_0.5.recall
\newcommand{\QOneRecall}{\ensuremath{76.8}\%}
% q1_results.json: main_with_min_filter.ks_yes_vs_no.D
\newcommand{\QOneKsD}{\ensuremath{0.673}}
% q1_results.json: main_with_min_filter.bins_reject_bh
\newcommand{\QOneBinsBh}{\ensuremath{5}}
% q1_results.json: main_with_min_filter.bins_reject_bonferroni
\newcommand{\QOneBinsBonferroni}{\ensuremath{4}}
% q1_results.json: no_min_filter.n_markets
\newcommand{\QOneNoMinN}{\ensuremath{74{,}474}}
% q1_results.json: no_min_filter.brier
\newcommand{\QOneNoMinBrier}{\ensuremath{0.134}}
% q1_results.json: no_min_filter.chi2_gof.chi2
\newcommand{\QOneNoMinChiSq}{\ensuremath{214.2}}
% q1_results.json: no_min_filter.chi2_gof.p
\newcommand{\QOneNoMinP}{\ensuremath{1.7\times10^{-40}}}
% q1_results.json: no_min_filter.chi2_gof.effect_w
\newcommand{\QOneNoMinW}{\ensuremath{0.054}}
% q1_results.json: min_trades
\newcommand{\MinTrades}{\ensuremath{10}}
% q1_bins.csv: bin=0 n
\newcommand{\BinZeroTenN}{\ensuremath{18{,}209}}
% q1_bins.csv: bin=0 mean_price
\newcommand{\BinZeroTenPrice}{\ensuremath{1.8}\%}
% q1_bins.csv: bin=0 observed_rate
\newcommand{\BinZeroTenObs}{\ensuremath{1.4}\%}
% q1_bins.csv: bin=0 diff_obs_minus_price
\newcommand{\BinZeroTenDiff}{\ensuremath{-0.46}\%}
% q1_bins.csv: bin=0 ci_low
\newcommand{\BinZeroTenCiLow}{\ensuremath{1.2}\%}
% q1_bins.csv: bin=0 ci_high
\newcommand{\BinZeroTenCiHigh}{\ensuremath{1.5}\%}
% q1_bins.csv: bin=0 binom_p
\newcommand{\BinZeroTenP}{\ensuremath{1.6\times10^{-6}}}
% q1_bins.csv: bin=0 p_bh
\newcommand{\BinZeroTenBhP}{\ensuremath{1.6\times10^{-5}}}
% q1_bins.csv: bin=1 n
\newcommand{\BinTenTwentyN}{\ensuremath{2{,}703}}
% q1_bins.csv: bin=1 mean_price
\newcommand{\BinTenTwentyPrice}{\ensuremath{14.8}\%}
% q1_bins.csv: bin=1 observed_rate
\newcommand{\BinTenTwentyObs}{\ensuremath{12.5}\%}
% q1_bins.csv: bin=1 diff_obs_minus_price
\newcommand{\BinTenTwentyDiff}{\ensuremath{-2.30}\%}
% q1_bins.csv: bin=1 ci_low
\newcommand{\BinTenTwentyCiLow}{\ensuremath{11.3}\%}
% q1_bins.csv: bin=1 ci_high
\newcommand{\BinTenTwentyCiHigh}{\ensuremath{13.8}\%}
% q1_bins.csv: bin=1 binom_p
\newcommand{\BinTenTwentyP}{\ensuremath{6.5\times10^{-4}}}
% q1_bins.csv: bin=1 p_bh
\newcommand{\BinTenTwentyBhP}{\ensuremath{0.002}}
% q1_bins.csv: bin=2 n
\newcommand{\BinTwentyThirtyN}{\ensuremath{2{,}543}}
% q1_bins.csv: bin=2 mean_price
\newcommand{\BinTwentyThirtyPrice}{\ensuremath{24.9}\%}
% q1_bins.csv: bin=2 observed_rate
\newcommand{\BinTwentyThirtyObs}{\ensuremath{23.9}\%}
% q1_bins.csv: bin=2 diff_obs_minus_price
\newcommand{\BinTwentyThirtyDiff}{\ensuremath{-1.00}\%}
% q1_bins.csv: bin=2 ci_low
\newcommand{\BinTwentyThirtyCiLow}{\ensuremath{22.3}\%}
% q1_bins.csv: bin=2 ci_high
\newcommand{\BinTwentyThirtyCiHigh}{\ensuremath{25.6}\%}
% q1_bins.csv: bin=2 binom_p
\newcommand{\BinTwentyThirtyP}{\ensuremath{0.252}}
% q1_bins.csv: bin=2 p_bh
\newcommand{\BinTwentyThirtyBhP}{\ensuremath{0.315}}
% q1_bins.csv: bin=3 n
\newcommand{\BinThirtyFortyN}{\ensuremath{2{,}769}}
% q1_bins.csv: bin=3 mean_price
\newcommand{\BinThirtyFortyPrice}{\ensuremath{34.9}\%}
% q1_bins.csv: bin=3 observed_rate
\newcommand{\BinThirtyFortyObs}{\ensuremath{33.2}\%}
% q1_bins.csv: bin=3 diff_obs_minus_price
\newcommand{\BinThirtyFortyDiff}{\ensuremath{-1.72}\%}
% q1_bins.csv: bin=3 ci_low
\newcommand{\BinThirtyFortyCiLow}{\ensuremath{31.5}\%}
% q1_bins.csv: bin=3 ci_high
\newcommand{\BinThirtyFortyCiHigh}{\ensuremath{35.0}\%}
% q1_bins.csv: bin=3 binom_p
\newcommand{\BinThirtyFortyP}{\ensuremath{0.058}}
% q1_bins.csv: bin=3 p_bh
\newcommand{\BinThirtyFortyBhP}{\ensuremath{0.097}}
% q1_bins.csv: bin=4 n
\newcommand{\BinFortyFiftyN}{\ensuremath{3{,}859}}
% q1_bins.csv: bin=4 mean_price
\newcommand{\BinFortyFiftyPrice}{\ensuremath{45.0}\%}
% q1_bins.csv: bin=4 observed_rate
\newcommand{\BinFortyFiftyObs}{\ensuremath{43.8}\%}
% q1_bins.csv: bin=4 diff_obs_minus_price
\newcommand{\BinFortyFiftyDiff}{\ensuremath{-1.23}\%}
% q1_bins.csv: bin=4 ci_low
\newcommand{\BinFortyFiftyCiLow}{\ensuremath{42.3}\%}
% q1_bins.csv: bin=4 ci_high
\newcommand{\BinFortyFiftyCiHigh}{\ensuremath{45.4}\%}
% q1_bins.csv: bin=4 binom_p
\newcommand{\BinFortyFiftyP}{\ensuremath{0.128}}
% q1_bins.csv: bin=4 p_bh
\newcommand{\BinFortyFiftyBhP}{\ensuremath{0.183}}
% q1_bins.csv: bin=5 n
\newcommand{\BinFiftySixtyN}{\ensuremath{4{,}240}}
% q1_bins.csv: bin=5 mean_price
\newcommand{\BinFiftySixtyPrice}{\ensuremath{54.1}\%}
% q1_bins.csv: bin=5 observed_rate
\newcommand{\BinFiftySixtyObs}{\ensuremath{53.5}\%}
% q1_bins.csv: bin=5 diff_obs_minus_price
\newcommand{\BinFiftySixtyDiff}{\ensuremath{-0.58}\%}
% q1_bins.csv: bin=5 ci_low
\newcommand{\BinFiftySixtyCiLow}{\ensuremath{52.0}\%}
% q1_bins.csv: bin=5 ci_high
\newcommand{\BinFiftySixtyCiHigh}{\ensuremath{55.0}\%}
% q1_bins.csv: bin=5 binom_p
\newcommand{\BinFiftySixtyP}{\ensuremath{0.450}}
% q1_bins.csv: bin=5 p_bh
\newcommand{\BinFiftySixtyBhP}{\ensuremath{0.500}}
% q1_bins.csv: bin=6 n
\newcommand{\BinSixtySeventyN}{\ensuremath{2{,}553}}
% q1_bins.csv: bin=6 mean_price
\newcommand{\BinSixtySeventyPrice}{\ensuremath{64.4}\%}
% q1_bins.csv: bin=6 observed_rate
\newcommand{\BinSixtySeventyObs}{\ensuremath{61.7}\%}
% q1_bins.csv: bin=6 diff_obs_minus_price
\newcommand{\BinSixtySeventyDiff}{\ensuremath{-2.79}\%}
% q1_bins.csv: bin=6 ci_low
\newcommand{\BinSixtySeventyCiLow}{\ensuremath{59.8}\%}
% q1_bins.csv: bin=6 ci_high
\newcommand{\BinSixtySeventyCiHigh}{\ensuremath{63.5}\%}
% q1_bins.csv: bin=6 binom_p
\newcommand{\BinSixtySeventyP}{\ensuremath{0.004}}
% q1_bins.csv: bin=6 p_bh
\newcommand{\BinSixtySeventyBhP}{\ensuremath{0.009}}
% q1_bins.csv: bin=7 n
\newcommand{\BinSeventyEightyN}{\ensuremath{2{,}301}}
% q1_bins.csv: bin=7 mean_price
\newcommand{\BinSeventyEightyPrice}{\ensuremath{74.6}\%}
% q1_bins.csv: bin=7 observed_rate
\newcommand{\BinSeventyEightyObs}{\ensuremath{71.3}\%}
% q1_bins.csv: bin=7 diff_obs_minus_price
\newcommand{\BinSeventyEightyDiff}{\ensuremath{-3.29}\%}
% q1_bins.csv: bin=7 ci_low
\newcommand{\BinSeventyEightyCiLow}{\ensuremath{69.4}\%}
% q1_bins.csv: bin=7 ci_high
\newcommand{\BinSeventyEightyCiHigh}{\ensuremath{73.1}\%}
% q1_bins.csv: bin=7 binom_p
\newcommand{\BinSeventyEightyP}{\ensuremath{3.6\times10^{-4}}}
% q1_bins.csv: bin=7 p_bh
\newcommand{\BinSeventyEightyBhP}{\ensuremath{0.002}}
% q1_bins.csv: bin=8 n
\newcommand{\BinEightyNinetyN}{\ensuremath{2{,}133}}
% q1_bins.csv: bin=8 mean_price
\newcommand{\BinEightyNinetyPrice}{\ensuremath{84.5}\%}
% q1_bins.csv: bin=8 observed_rate
\newcommand{\BinEightyNinetyObs}{\ensuremath{82.7}\%}
% q1_bins.csv: bin=8 diff_obs_minus_price
\newcommand{\BinEightyNinetyDiff}{\ensuremath{-1.88}\%}
% q1_bins.csv: bin=8 ci_low
\newcommand{\BinEightyNinetyCiLow}{\ensuremath{81.0}\%}
% q1_bins.csv: bin=8 ci_high
\newcommand{\BinEightyNinetyCiHigh}{\ensuremath{84.2}\%}
% q1_bins.csv: bin=8 binom_p
\newcommand{\BinEightyNinetyP}{\ensuremath{0.018}}
% q1_bins.csv: bin=8 p_bh
\newcommand{\BinEightyNinetyBhP}{\ensuremath{0.036}}
% q1_bins.csv: bin=9 n
\newcommand{\BinNinetyHundredN}{\ensuremath{6{,}638}}
% q1_bins.csv: bin=9 mean_price
\newcommand{\BinNinetyHundredPrice}{\ensuremath{97.7}\%}
% q1_bins.csv: bin=9 observed_rate
\newcommand{\BinNinetyHundredObs}{\ensuremath{97.8}\%}
% q1_bins.csv: bin=9 diff_obs_minus_price
\newcommand{\BinNinetyHundredDiff}{\ensuremath{0.03}\%}
% q1_bins.csv: bin=9 ci_low
\newcommand{\BinNinetyHundredCiLow}{\ensuremath{97.4}\%}
% q1_bins.csv: bin=9 ci_high
\newcommand{\BinNinetyHundredCiHigh}{\ensuremath{98.1}\%}
% q1_bins.csv: bin=9 binom_p
\newcommand{\BinNinetyHundredP}{\ensuremath{0.902}}
% q1_bins.csv: bin=9 p_bh
\newcommand{\BinNinetyHundredBhP}{\ensuremath{0.902}}
% q2_results.json: concentration.maker.n_wallets
\newcommand{\MakerNWallets}{\ensuremath{581{,}435}}
% q2_results.json: concentration.maker.gini
\newcommand{\MakerGini}{\ensuremath{0.978}}
% q2_results.json: concentration.maker.top1_share
\newcommand{\MakerTopOneShare}{\ensuremath{86.1}\%}
% q2_results.json: concentration.maker.top5_share
\newcommand{\MakerTopFiveShare}{\ensuremath{94.6}\%}
% q2_results.json: concentration.maker.top10_share
\newcommand{\MakerTopTenShare}{\ensuremath{96.9}\%}
% q2_results.json: concentration.maker.mean
\newcommand{\MakerMeanVolume}{\ensuremath{14{,}103}}
% q2_results.json: concentration.maker.median
\newcommand{\MakerMedianVolume}{\ensuremath{101.2}}
% q2_results.json: concentration.taker.n_wallets
\newcommand{\TakerNWallets}{\ensuremath{1{,}535{,}977}}
% q2_results.json: concentration.taker.gini
\newcommand{\TakerGini}{\ensuremath{0.963}}
% q2_results.json: concentration.taker.top1_share
\newcommand{\TakerTopOneShare}{\ensuremath{75.9}\%}
% q2_results.json: concentration.taker.top5_share
\newcommand{\TakerTopFiveShare}{\ensuremath{90.4}\%}
% q2_results.json: concentration.taker.top10_share
\newcommand{\TakerTopTenShare}{\ensuremath{94.8}\%}
% q2_results.json: concentration.taker.mean
\newcommand{\TakerMeanVolume}{\ensuremath{5{,}339}}
% q2_results.json: concentration.taker.median
\newcommand{\TakerMedianVolume}{\ensuremath{86.6}}
% q2_results.json: role_tests.welch_t.t
\newcommand{\RoleWelchT}{\ensuremath{117.1}}
% q2_results.json: role_tests.welch_t.p
\newcommand{\RoleWelchP}{\ensuremath{<10^{-300}}}
% q2_results.json: role_tests.welch_t.cohens_d
\newcommand{\RoleCohenD}{\ensuremath{0.177}}
% q2_results.json: role_tests.ks.D
\newcommand{\RoleKsD}{\ensuremath{0.057}}
% q2_results.json: spearman_n_trades_vs_volume.rho
\newcommand{\SpearmanTradesVolumeRho}{\ensuremath{0.655}}
% q2_results.json: kmeans.anova_mean_trade_price.eta_squared
\newcommand{\KmeansPriceEtaSq}{\ensuremath{0.008}}
% q2_results.json: kmeans.anova_mean_trade_price.perm_p
\newcommand{\KmeansPricePermP}{\ensuremath{1.0\times10^{-3}}}
% q2_results.json: kmeans.k_chosen
\newcommand{\KmeansK}{\ensuremath{4}}
% q2_results.json: kmeans.n_sample
\newcommand{\KmeansNSample}{\ensuremath{200{,}000}}
% q2_results.json: kmeans.robust_scaling_check.adjusted_rand_index_vs_zscore
\newcommand{\KmeansAri}{\ensuremath{0.928}}
% q2_results.json: kmeans.anova_mean_trade_price.n_perm
\newcommand{\KmeansNPerm}{\ensuremath{1{,}000}}
% q2_results.json: kmeans.chi2_cluster_x_category.cramers_v
\newcommand{\KmeansCategoryV}{\ensuremath{0.105}}
% q2_results.json: kmeans.chi2_cluster_x_category.chi2
\newcommand{\KmeansCategoryChi}{\ensuremath{6{,}653.0}}
% q2_results.json: kmeans.chi2_cluster_x_category.df
\newcommand{\KmeansCategoryDf}{\ensuremath{21}}
% q2_results.json: kmeans.chi2_cluster_x_category.p
\newcommand{\KmeansCategoryP}{\ensuremath{<10^{-300}}}
% q2_results.json: kmeans.cluster_summary[cluster=0].n
\newcommand{\ClusterZeroN}{\ensuremath{57{,}550}}
% q2_results.json: kmeans.cluster_summary[cluster=0].median_volume
\newcommand{\ClusterZeroMedianVolume}{\ensuremath{212.7}}
% q2_results.json: kmeans.cluster_summary[cluster=0].median_n_trades
\newcommand{\ClusterZeroMedianTrades}{\ensuremath{6}}
% q2_results.json: kmeans.cluster_summary[cluster=0].median_n_markets
\newcommand{\ClusterZeroMedianMarkets}{\ensuremath{2}}
% q2_results.json: kmeans.cluster_summary[cluster=1].n
\newcommand{\ClusterOneN}{\ensuremath{15{,}294}}
% q2_results.json: kmeans.cluster_summary[cluster=1].median_volume
\newcommand{\ClusterOneMedianVolume}{\ensuremath{5{,}815.3}}
% q2_results.json: kmeans.cluster_summary[cluster=1].median_n_trades
\newcommand{\ClusterOneMedianTrades}{\ensuremath{129}}
% q2_results.json: kmeans.cluster_summary[cluster=1].median_n_markets
\newcommand{\ClusterOneMedianMarkets}{\ensuremath{42}}
% q2_results.json: kmeans.cluster_summary[cluster=2].n
\newcommand{\ClusterTwoN}{\ensuremath{69{,}556}}
% q2_results.json: kmeans.cluster_summary[cluster=2].median_volume
\newcommand{\ClusterTwoMedianVolume}{\ensuremath{6.6}}
% q2_results.json: kmeans.cluster_summary[cluster=2].median_n_trades
\newcommand{\ClusterTwoMedianTrades}{\ensuremath{2}}
% q2_results.json: kmeans.cluster_summary[cluster=2].median_n_markets
\newcommand{\ClusterTwoMedianMarkets}{\ensuremath{1}}
% q2_results.json: kmeans.cluster_summary[cluster=3].n
\newcommand{\ClusterThreeN}{\ensuremath{57{,}600}}
% q2_results.json: kmeans.cluster_summary[cluster=3].median_volume
\newcommand{\ClusterThreeMedianVolume}{\ensuremath{269.0}}
% q2_results.json: kmeans.cluster_summary[cluster=3].median_n_trades
\newcommand{\ClusterThreeMedianTrades}{\ensuremath{20}}
% q2_results.json: kmeans.cluster_summary[cluster=3].median_n_markets
\newcommand{\ClusterThreeMedianMarkets}{\ensuremath{10}}
% q3_results.json: summary.p_1d.mean
\newcommand{\PriceOneDayMean}{\ensuremath{0.376}}
% q3_results.json: summary.p_1d.median
\newcommand{\PriceOneDayMedian}{\ensuremath{0.320}}
% q3_results.json: summary.trade_size_usdc.mean
\newcommand{\TradeSizeMean}{\ensuremath{129.8}}
% q3_results.json: summary.trade_size_usdc.median
\newcommand{\TradeSizeMedian}{\ensuremath{6.9}}
% q3_results.json: fits.trade_size_usdc.power_law.alpha
\newcommand{\TradeSizePowerAlpha}{\ensuremath{1.875}}
% q3_results.json: fits.trade_size_usdc.power_law.xmin
\newcommand{\TradeSizePowerXmin}{\ensuremath{154.0}}
% q3_results.json: fits.trade_size_usdc.power_law.D
\newcommand{\TradeSizePowerKsD}{\ensuremath{0.039}}
% q3_results.json: fits.trade_size_usdc.exponential.D
\newcommand{\TradeSizeExponentialKsD}{\ensuremath{0.534}}
% q3_results.json: fits.trade_size_usdc.normal.D
\newcommand{\TradeSizeNormalKsD}{\ensuremath{0.469}}
% q3_results.json: fits.trade_size_usdc.power_law.p
\newcommand{\TradeSizePowerKsP}{\ensuremath{1.5\times10^{-131}}}
% q3_results.json: fits.trade_size_usdc.power_law.n_tail
\newcommand{\TradeSizePowerNTail}{\ensuremath{100{,}090}}
% q3_results.json: fits.wallet_volume_usdc.power_law.alpha
\newcommand{\WalletVolumePowerAlpha}{\ensuremath{1.788}}
% q3_results.json: fits.wallet_volume_usdc.power_law.xmin
\newcommand{\WalletVolumePowerXmin}{\ensuremath{3{,}450.2}}
% q3_results.json: fits.wallet_volume_usdc.power_law.D
\newcommand{\WalletVolumePowerKsD}{\ensuremath{0.024}}
% q3_results.json: fits.wallet_volume_usdc.exponential.D
\newcommand{\WalletVolumeExponentialKsD}{\ensuremath{0.675}}
% q3_results.json: fits.wallet_volume_usdc.normal.D
\newcommand{\WalletVolumeNormalKsD}{\ensuremath{0.488}}
% q3_results.json: fits.wallet_volume_usdc.power_law.p
\newcommand{\WalletVolumePowerKsP}{\ensuremath{1.2\times10^{-105}}}
% q3_results.json: fits.wallet_volume_usdc.power_law.n_tail
\newcommand{\WalletVolumePowerNTail}{\ensuremath{211{,}742}}
% lecture_numbers.json: z_interval_95.n_outside
\newcommand{\ZIntervalBinsOutsideN}{\ensuremath{5}}
% lecture_numbers.json: z_interval_95.bins_outside
\newcommand{\ZIntervalBinsOutsideList}{0, 1, 6, 7, 8}
% lecture_numbers.json: z_interval_95.bins_bh_off
\newcommand{\ZIntervalBinsBhOffList}{0, 1, 6, 7, 8}
% lecture_numbers.json: z_interval_95.same_bins
\newcommand{\ZIntervalSameAsBh}{yes}
% lecture_numbers.json: z_interval_95.bins_equal_frequency[price_outside]
\newcommand{\ZIntervalEqFreqOutsideN}{\ensuremath{8}}
% lecture_numbers.json: z_interval_95.bins[bin=0].bernoulli_sd
\newcommand{\BinZeroTenZSd}{\ensuremath{0.116}}
% lecture_numbers.json: z_interval_95.bins[bin=0].z_low
\newcommand{\BinZeroTenZLow}{\ensuremath{1.2}\%}
% lecture_numbers.json: z_interval_95.bins[bin=0].z_high
\newcommand{\BinZeroTenZHigh}{\ensuremath{1.5}\%}
% lecture_numbers.json: z_interval_95.bins[bin=0].price_outside
\newcommand{\BinZeroTenZOutside}{yes}
% lecture_numbers.json: z_interval_95.bins[bin=1].bernoulli_sd
\newcommand{\BinTenTwentyZSd}{\ensuremath{0.331}}
% lecture_numbers.json: z_interval_95.bins[bin=1].z_low
\newcommand{\BinTenTwentyZLow}{\ensuremath{11.3}\%}
% lecture_numbers.json: z_interval_95.bins[bin=1].z_high
\newcommand{\BinTenTwentyZHigh}{\ensuremath{13.8}\%}
% lecture_numbers.json: z_interval_95.bins[bin=1].price_outside
\newcommand{\BinTenTwentyZOutside}{yes}
% lecture_numbers.json: z_interval_95.bins[bin=2].bernoulli_sd
\newcommand{\BinTwentyThirtyZSd}{\ensuremath{0.427}}
% lecture_numbers.json: z_interval_95.bins[bin=2].z_low
\newcommand{\BinTwentyThirtyZLow}{\ensuremath{22.3}\%}
% lecture_numbers.json: z_interval_95.bins[bin=2].z_high
\newcommand{\BinTwentyThirtyZHigh}{\ensuremath{25.6}\%}
% lecture_numbers.json: z_interval_95.bins[bin=2].price_outside
\newcommand{\BinTwentyThirtyZOutside}{no}
% lecture_numbers.json: z_interval_95.bins[bin=3].bernoulli_sd
\newcommand{\BinThirtyFortyZSd}{\ensuremath{0.471}}
% lecture_numbers.json: z_interval_95.bins[bin=3].z_low
\newcommand{\BinThirtyFortyZLow}{\ensuremath{31.4}\%}
% lecture_numbers.json: z_interval_95.bins[bin=3].z_high
\newcommand{\BinThirtyFortyZHigh}{\ensuremath{34.9}\%}
% lecture_numbers.json: z_interval_95.bins[bin=3].price_outside
\newcommand{\BinThirtyFortyZOutside}{no}
% lecture_numbers.json: z_interval_95.bins[bin=4].bernoulli_sd
\newcommand{\BinFortyFiftyZSd}{\ensuremath{0.496}}
% lecture_numbers.json: z_interval_95.bins[bin=4].z_low
\newcommand{\BinFortyFiftyZLow}{\ensuremath{42.3}\%}
% lecture_numbers.json: z_interval_95.bins[bin=4].z_high
\newcommand{\BinFortyFiftyZHigh}{\ensuremath{45.4}\%}
% lecture_numbers.json: z_interval_95.bins[bin=4].price_outside
\newcommand{\BinFortyFiftyZOutside}{no}
% lecture_numbers.json: z_interval_95.bins[bin=5].bernoulli_sd
\newcommand{\BinFiftySixtyZSd}{\ensuremath{0.499}}
% lecture_numbers.json: z_interval_95.bins[bin=5].z_low
\newcommand{\BinFiftySixtyZLow}{\ensuremath{52.0}\%}
% lecture_numbers.json: z_interval_95.bins[bin=5].z_high
\newcommand{\BinFiftySixtyZHigh}{\ensuremath{55.0}\%}
% lecture_numbers.json: z_interval_95.bins[bin=5].price_outside
\newcommand{\BinFiftySixtyZOutside}{no}
% lecture_numbers.json: z_interval_95.bins[bin=6].bernoulli_sd
\newcommand{\BinSixtySeventyZSd}{\ensuremath{0.486}}
% lecture_numbers.json: z_interval_95.bins[bin=6].z_low
\newcommand{\BinSixtySeventyZLow}{\ensuremath{59.8}\%}
% lecture_numbers.json: z_interval_95.bins[bin=6].z_high
\newcommand{\BinSixtySeventyZHigh}{\ensuremath{63.5}\%}
% lecture_numbers.json: z_interval_95.bins[bin=6].price_outside
\newcommand{\BinSixtySeventyZOutside}{yes}
% lecture_numbers.json: z_interval_95.bins[bin=7].bernoulli_sd
\newcommand{\BinSeventyEightyZSd}{\ensuremath{0.452}}
% lecture_numbers.json: z_interval_95.bins[bin=7].z_low
\newcommand{\BinSeventyEightyZLow}{\ensuremath{69.4}\%}
% lecture_numbers.json: z_interval_95.bins[bin=7].z_high
\newcommand{\BinSeventyEightyZHigh}{\ensuremath{73.1}\%}
% lecture_numbers.json: z_interval_95.bins[bin=7].price_outside
\newcommand{\BinSeventyEightyZOutside}{yes}
% lecture_numbers.json: z_interval_95.bins[bin=8].bernoulli_sd
\newcommand{\BinEightyNinetyZSd}{\ensuremath{0.379}}
% lecture_numbers.json: z_interval_95.bins[bin=8].z_low
\newcommand{\BinEightyNinetyZLow}{\ensuremath{81.0}\%}
% lecture_numbers.json: z_interval_95.bins[bin=8].z_high
\newcommand{\BinEightyNinetyZHigh}{\ensuremath{84.3}\%}
% lecture_numbers.json: z_interval_95.bins[bin=8].price_outside
\newcommand{\BinEightyNinetyZOutside}{yes}
% lecture_numbers.json: z_interval_95.bins[bin=9].bernoulli_sd
\newcommand{\BinNinetyHundredZSd}{\ensuremath{0.148}}
% lecture_numbers.json: z_interval_95.bins[bin=9].z_low
\newcommand{\BinNinetyHundredZLow}{\ensuremath{97.4}\%}
% lecture_numbers.json: z_interval_95.bins[bin=9].z_high
\newcommand{\BinNinetyHundredZHigh}{\ensuremath{98.1}\%}
% lecture_numbers.json: z_interval_95.bins[bin=9].price_outside
\newcommand{\BinNinetyHundredZOutside}{no}
% lecture_numbers.json: largest_gap.bin
\newcommand{\LargestGapBin}{\ensuremath{7}}
% lecture_numbers.json: largest_gap.diff_obs_minus_price
\newcommand{\LargestGapDiff}{\ensuremath{-3.29}\%}
% lecture_numbers.json: largest_gap.abs
\newcommand{\LargestGapAbs}{\ensuremath{3.29}\%}
% lecture_numbers.json: log10_mean_difference_maker_minus_taker.difference
\newcommand{\RoleLogTenDiff}{\ensuremath{0.219}}
% lecture_numbers.json: log10_mean_difference_maker_minus_taker.mean_log10_maker
\newcommand{\RoleLogTenMaker}{\ensuremath{2.104}}
% lecture_numbers.json: log10_mean_difference_maker_minus_taker.mean_log10_taker
\newcommand{\RoleLogTenTaker}{\ensuremath{1.885}}
% lecture_numbers.json: top20_share.maker
\newcommand{\MakerTopTwentyShare}{\ensuremath{98.9}\%}
% lecture_numbers.json: top20_share.taker
\newcommand{\TakerTopTwentyShare}{\ensuremath{98.0}\%}
% lecture_numbers.json: top20_share.reference
\newcommand{\TopTwentyReference}{\ensuremath{80.0}\%}
% lecture_numbers.json: gap_direction.main.n_negative
\newcommand{\GapNegBinsMain}{\ensuremath{9}}
% lecture_numbers.json: gap_direction.main.net_gap (n-weighted)
\newcommand{\GapNetDiffMain}{\ensuremath{-0.99}\%}
% lecture_numbers.json: gap_direction.no_min_filter.n_negative
\newcommand{\GapNegBinsNoMin}{\ensuremath{9}}
% lecture_numbers.json: gap_direction.no_min_filter.net_gap (n-weighted)
\newcommand{\GapNetDiffNoMin}{\ensuremath{-1.03}\%}
% lecture_numbers.json: gap_direction.equal_frequency.n_negative
\newcommand{\GapNegBinsEqFreq}{\ensuremath{9}}
% lecture_numbers.json: gap_direction.equal_frequency.net_gap (n-weighted)
\newcommand{\GapNetDiffEqFreq}{\ensuremath{-0.99}\%}
% lecture_numbers.json: gap_direction.main.n_bins
\newcommand{\GapBinsTotal}{\ensuremath{10}}
% lecture_numbers.json: gap_direction.main.n_yes_minus_expected (absolute value)
\newcommand{\GapNetYesShort}{\ensuremath{475.2}}
% lecture_numbers.json: chi2_largest_contributor.bin
\newcommand{\ChiTopBin}{\ensuremath{0}}
% lecture_numbers.json: chi2_largest_contributor.contribution
\newcommand{\ChiTopContribution}{\ensuremath{21.3}}
% lecture_numbers.json: chi2_largest_contributor.share
\newcommand{\ChiTopShare}{\ensuremath{31.3}\%}
% lecture_numbers.json: no_min_filter_longshot_bh.0.diff_obs_minus_price
\newcommand{\NoMinBinZeroDiff}{\ensuremath{0.05}\%}
% lecture_numbers.json: no_min_filter_longshot_bh.0.p_bh
\newcommand{\NoMinBinZeroPBh}{\ensuremath{0.601}}
% lecture_numbers.json: no_min_filter_longshot_bh.0.reject_bh
\newcommand{\NoMinBinZeroBhOff}{no}
% lecture_numbers.json: no_min_filter_longshot_bh.1.diff_obs_minus_price
\newcommand{\NoMinBinOneDiff}{\ensuremath{-0.93}\%}
% lecture_numbers.json: no_min_filter_longshot_bh.1.p_bh
\newcommand{\NoMinBinOnePBh}{\ensuremath{0.136}}
% lecture_numbers.json: no_min_filter_longshot_bh.1.reject_bh
\newcommand{\NoMinBinOneBhOff}{no}
% lecture_numbers.json: cluster_category_gap.cluster
\newcommand{\ClusterCategoryGapCluster}{\ensuremath{2}}
% lecture_numbers.json: cluster_category_gap.category
\newcommand{\ClusterCategoryGapCategory}{Sports}
% lecture_numbers.json: cluster_category_gap.observed_share
\newcommand{\ClusterCategoryGapObs}{\ensuremath{6.7}\%}
% lecture_numbers.json: cluster_category_gap.expected_share
\newcommand{\ClusterCategoryGapExp}{\ensuremath{8.2}\%}
% lecture_numbers.json: cluster_category_gap.diff
\newcommand{\ClusterCategoryGapDiff}{\ensuremath{-1.49}\%}
% lecture_numbers.json: ssb_ratio.mean_trade_price
\newcommand{\KmeansPriceSsbRatio}{\ensuremath{0.008}}
% lecture_numbers.json: ssb_ratio.log_median_trade_size
\newcommand{\KmeansTradeSizeSsbRatio}{\ensuremath{0.224}}
% label1_counts.json (run 150296_v1label): overall.share_yes_like
\newcommand{\OutcomeOneYesShare}{\ensuremath{63.9}\%}
% label1_counts.json (run 150296_v1label): overall.n_yes_like
\newcommand{\OutcomeOneYesN}{\ensuremath{30{,}630}}
% label1_counts.json (run 150296_v1label): by_bin_yes_only bins 0 to 8 with diff_obs_minus_price < 0
\newcommand{\OutcomeOneYesBinsNegative}{\ensuremath{9}}
% label1_counts.json (run 150296_v1label): by_bin_yes_only bins 0 to 8, smallest |diff_obs_minus_price|, percentage points
\newcommand{\OutcomeOneYesGapMin}{\ensuremath{0.4}}
% label1_counts.json (run 150296_v1label): by_bin_yes_only bins 0 to 8, largest |diff_obs_minus_price|, percentage points
\newcommand{\OutcomeOneYesGapMax}{\ensuremath{5.6}}
% q1_bins_eqfreq.csv: row count
\newcommand{\QOneEqFreqBins}{\ensuremath{10}}
% 142140_q1eqfreq_v5/q1_bins_eqfreq.csv: reject_bh (saved run output, code v5)
\newcommand{\QOneEqFreqBinsBh}{\ensuremath{7}}
% close_time_summary.json (pilot run): n_markets
\newcommand{\PilotNMarkets}{\ensuremath{2{,}063}}
% close_time_summary.json (pilot run): share_markets_trade_after_close
\newcommand{\PilotShareMarketsAfterClose}{\ensuremath{51.3}\%}
% close_time_summary.json (pilot run): share_volume_after_close
\newcommand{\PilotShareVolumeAfterClose}{\ensuremath{34.0}\%}
% robustness run directory name
\newcommand{\RobustRun}{142123\_robust\_full\_v4}
% manifest_r1_split_payout.json: code_sha (hash of the script r1_split_payout.py alone)
\newcommand{\RobustCodeShaROne}{1232fa0fe1e2}
% manifest_r2_price_path.json: code_sha (hash of the script r2_price_path.py alone)
\newcommand{\RobustCodeShaRTwo}{4bcaceda3b5a}
% manifest_r1_split_payout.json, manifest_r2_price_path.json: code_version
\newcommand{\RobustCodeVersion}{v4}
% code.sha256: sha256 of the file that lists the hash of every script (first 12 hex)
\newcommand{\RobustCodeSha}{f91fc1970be1}
% path of the per-script hash list
\newcommand{\RobustCodeShaFile}{runs/142123\_robust\_full\_v4/code.sha256}
% r1_results.json: n_kept
\newcommand{\RobROneNKept}{\ensuremath{47{,}948}}
% r1_results.json: class_counts.clean
\newcommand{\RobROneCleanN}{\ensuremath{47{,}940}}
% r1_results.json: class_shares.clean
\newcommand{\RobROneCleanShare}{\ensuremath{99.98}\%}
% r1_results.json: class_counts.split
\newcommand{\RobROneSplitN}{\ensuremath{7}}
% r1_results.json: class_shares.split
\newcommand{\RobROneSplitShare}{\ensuremath{0.01}\%}
% r1_results.json: class_counts.other
\newcommand{\RobROneOtherN}{\ensuremath{0}}
% r1_results.json: class_shares.other
\newcommand{\RobROneOtherShare}{\ensuremath{0.00}\%}
% r1_results.json: class_counts.missing
\newcommand{\RobROneMissingN}{\ensuremath{1}}
% r1_results.json: class_shares.missing
\newcommand{\RobROneMissingShare}{\ensuremath{0.00}\%}
% r1_results.json: duplicates.kept_ids_with_duplicate_rows
\newcommand{\RobROneDupIds}{\ensuremath{0}}
% r1_results.json: duplicates.kept_duplicates_conflicting
\newcommand{\RobROneDupConflicting}{\ensuremath{0}}
% r1_results.json: our_y_on_split_markets.y1
\newcommand{\RobROneSplitYesOne}{\ensuremath{7}}
% r1_results.json: our_y_on_split_markets.y0
\newcommand{\RobROneSplitYesZero}{\ensuremath{0}}
% r1_results.json: mapping.n_clean
\newcommand{\RobROneNClean}{\ensuremath{47{,}940}}
% r1_results.json: mapping.agreement_index_primary
\newcommand{\RobROneAgreement}{\ensuremath{99.92}\%}
% r1_results.json: mapping.agreement_ci95[0]
\newcommand{\RobROneAgreementLow}{\ensuremath{99.89}\%}
% r1_results.json: mapping.agreement_ci95[1]
\newcommand{\RobROneAgreementHigh}{\ensuremath{99.94}\%}
% r1_results.json: mapping.n_mismatch
\newcommand{\RobROneNMismatch}{\ensuremath{40}}
% r1_results.json: mapping.validated
\newcommand{\RobROneMappingValidated}{yes}
% r1_results.json: sensitivity.n_excluded
\newcommand{\RobROneNExcluded}{\ensuremath{47}}
% r1_results.json: sensitivity.n_markets
\newcommand{\RobROneSensN}{\ensuremath{47{,}901}}
% r1_results.json: sensitivity.q1_sensitivity.chi2_gof.chi2
\newcommand{\RobROneSensChiSq}{\ensuremath{64.3}}
% r1_results.json: sensitivity.q1_sensitivity.chi2_gof.df
\newcommand{\RobROneSensDf}{\ensuremath{10}}
% r1_results.json: sensitivity.q1_sensitivity.chi2_gof.p
\newcommand{\RobROneSensP}{\ensuremath{5.6\times10^{-10}}}
% r1_results.json: sensitivity.q1_sensitivity.chi2_gof.effect_w
\newcommand{\RobROneSensW}{\ensuremath{0.037}}
% r1_results.json: sensitivity.q1_sensitivity.brier
\newcommand{\RobROneSensBrier}{\ensuremath{0.107}}
% r1_results.json: sensitivity.q1_sensitivity.bins_reject_bh
\newcommand{\RobROneSensBinsBh}{\ensuremath{4}}
% r1_results.json: sensitivity.delta_effect_w
\newcommand{\RobROneDeltaW}{\ensuremath{-0.001}}
% r1_results.json: sensitivity.rule_results_same
\newcommand{\RobROneRuleSame}{yes}
% r1_results.json: sensitivity.rule_result_original
\newcommand{\RobROneRuleOriginal}{Prices are not calibrated as a whole at this n.}
% r1_results.json: sensitivity.rule_result_sensitivity
\newcommand{\RobROneRuleSens}{Prices are not calibrated as a whole at this n.}
% r2_results.json: overall.n
\newcommand{\RobRTwoN}{\ensuremath{47{,}948}}
% r2_results.json: overall.median_abs_delta
\newcommand{\RobRTwoMedianAbsDelta}{\ensuremath{0.015}}
% r2_results.json: overall.share_no_trade
\newcommand{\RobRTwoShareNoTrade}{\ensuremath{5.9}\%}
% r2_results.json: overall.share_crossing
\newcommand{\RobRTwoShareCrossing}{\ensuremath{8.02}\%}
% r2_results.json: overall.share_abs_delta_ge_0.10
\newcommand{\RobRTwoShareDeltaAtLeastTen}{\ensuremath{18.98}\%}
% r2_results.json: overall.share_abs_delta_ge_0.25
\newcommand{\RobRTwoShareDeltaAtLeastTwentyFive}{\ensuremath{10.97}\%}
% r2_results.json: overall.share_abs_delta_ge_0.50
\newcommand{\RobRTwoShareDeltaAtLeastFifty}{\ensuremath{4.65}\%}
% r2_results.json: overall.quantiles_abs_delta.q0.5
\newcommand{\RobRTwoAbsDeltaMedian}{\ensuremath{0.015}}
% r2_results.json: overall.quantiles_abs_delta.q0.75
\newcommand{\RobRTwoAbsDeltaSevenFive}{\ensuremath{0.060}}
% r2_results.json: overall.quantiles_abs_delta.q0.9
\newcommand{\RobRTwoAbsDeltaNinety}{\ensuremath{0.280}}
% r2_results.json: overall.quantiles_abs_delta.q0.95
\newcommand{\RobRTwoAbsDeltaNinetyFive}{\ensuremath{0.482}}
% r2_results.json: overall.quantiles_abs_delta.q0.99
\newcommand{\RobRTwoAbsDeltaNinetyNine}{\ensuremath{0.789}}
% r2_results.json: overall.quantiles_abs_delta.max
\newcommand{\RobRTwoAbsDeltaMax}{\ensuremath{0.998}}
% r2_results.json: overall.n_window_trades.median
\newcommand{\RobRTwoWindowTradesMedian}{\ensuremath{61}}
% r2_results.json: overall.n_window_trades.q0.9
\newcommand{\RobRTwoWindowTradesNinety}{\ensuremath{1{,}147}}
% r2_results.json: overall.n_window_trades.max
\newcommand{\RobRTwoWindowTradesMax}{\ensuremath{40{,}499}}
% r2_results.json: overall.brier.n
\newcommand{\RobRTwoAllN}{\ensuremath{47{,}948}}
% r2_results.json: overall.brier.brier_1d
\newcommand{\RobRTwoAllBrierOneDay}{\ensuremath{0.107}}
% r2_results.json: overall.brier.brier_T
\newcommand{\RobRTwoAllBrierAtClose}{\ensuremath{0.075}}
% r2_results.json: overall.brier.mean_diff
\newcommand{\RobRTwoAllMeanDiff}{\ensuremath{-0.032}}
% r2_results.json: overall.brier.ci_low
\newcommand{\RobRTwoAllCiLow}{\ensuremath{-0.033}}
% r2_results.json: overall.brier.ci_high
\newcommand{\RobRTwoAllCiHigh}{\ensuremath{-0.031}}
% r2_results.json: overall.brier.t
\newcommand{\RobRTwoAllT}{\ensuremath{-56.5}}
% r2_results.json: overall.brier.p
\newcommand{\RobRTwoAllP}{\ensuremath{<10^{-300}}}
% r2_results.json: overall.brier.cohens_dz
\newcommand{\RobRTwoAllDz}{\ensuremath{-0.258}}
% r2_results.json: overall.brier_traded_only.n
\newcommand{\RobRTwoTradedN}{\ensuremath{45{,}112}}
% r2_results.json: overall.brier_traded_only.brier_1d
\newcommand{\RobRTwoTradedBrierOneDay}{\ensuremath{0.112}}
% r2_results.json: overall.brier_traded_only.brier_T
\newcommand{\RobRTwoTradedBrierAtClose}{\ensuremath{0.078}}
% r2_results.json: overall.brier_traded_only.mean_diff
\newcommand{\RobRTwoTradedMeanDiff}{\ensuremath{-0.034}}
% r2_results.json: overall.brier_traded_only.ci_low
\newcommand{\RobRTwoTradedCiLow}{\ensuremath{-0.035}}
% r2_results.json: overall.brier_traded_only.ci_high
\newcommand{\RobRTwoTradedCiHigh}{\ensuremath{-0.033}}
% r2_results.json: overall.brier_traded_only.t
\newcommand{\RobRTwoTradedT}{\ensuremath{-56.6}}
% r2_results.json: overall.brier_traded_only.p
\newcommand{\RobRTwoTradedP}{\ensuremath{<10^{-300}}}
% r2_results.json: overall.brier_traded_only.cohens_dz
\newcommand{\RobRTwoTradedDz}{\ensuremath{-0.266}}
% r2_results.json: by_category.Crypto.n
\newcommand{\RobRTwoCatCryptoN}{\ensuremath{11{,}888}}
% r2_results.json: by_category.Crypto.median_abs_delta
\newcommand{\RobRTwoCatCryptoMedianAbsDelta}{\ensuremath{0.004}}
% r2_results.json: by_category.Crypto.brier.mean_diff
\newcommand{\RobRTwoCatCryptoMeanDiff}{\ensuremath{-0.032}}
% r2_results.json: by_category.Culture.n
\newcommand{\RobRTwoCatCultureN}{\ensuremath{1{,}074}}
% r2_results.json: by_category.Culture.median_abs_delta
\newcommand{\RobRTwoCatCultureMedianAbsDelta}{\ensuremath{0.010}}
% r2_results.json: by_category.Culture.brier.mean_diff
\newcommand{\RobRTwoCatCultureMeanDiff}{\ensuremath{-0.008}}
% r2_results.json: by_category.Finance.n
\newcommand{\RobRTwoCatFinanceN}{\ensuremath{5{,}146}}
% r2_results.json: by_category.Finance.median_abs_delta
\newcommand{\RobRTwoCatFinanceMedianAbsDelta}{\ensuremath{0.015}}
% r2_results.json: by_category.Finance.brier.mean_diff
\newcommand{\RobRTwoCatFinanceMeanDiff}{\ensuremath{-0.042}}
% r2_results.json: by_category.Other.n
\newcommand{\RobRTwoCatOtherN}{\ensuremath{2{,}274}}
% r2_results.json: by_category.Other.median_abs_delta
\newcommand{\RobRTwoCatOtherMedianAbsDelta}{\ensuremath{0.011}}
% r2_results.json: by_category.Other.brier.mean_diff
\newcommand{\RobRTwoCatOtherMeanDiff}{\ensuremath{-0.024}}
% r2_results.json: by_category.Politics.n
\newcommand{\RobRTwoCatPoliticsN}{\ensuremath{6{,}922}}
% r2_results.json: by_category.Politics.median_abs_delta
\newcommand{\RobRTwoCatPoliticsMedianAbsDelta}{\ensuremath{0.013}}
% r2_results.json: by_category.Politics.brier.mean_diff
\newcommand{\RobRTwoCatPoliticsMeanDiff}{\ensuremath{-0.014}}
% r2_results.json: by_category.Price Action.n
\newcommand{\RobRTwoCatPriceActionN}{\ensuremath{663}}
% r2_results.json: by_category.Price Action.median_abs_delta
\newcommand{\RobRTwoCatPriceActionMedianAbsDelta}{\ensuremath{0.003}}
% r2_results.json: by_category.Price Action.brier.mean_diff
\newcommand{\RobRTwoCatPriceActionMeanDiff}{\ensuremath{-0.047}}
% r2_results.json: by_category.Sci-Tech.n
\newcommand{\RobRTwoCatSciTechN}{\ensuremath{819}}
% r2_results.json: by_category.Sci-Tech.median_abs_delta
\newcommand{\RobRTwoCatSciTechMedianAbsDelta}{\ensuremath{0.007}}
% r2_results.json: by_category.Sci-Tech.brier.mean_diff
\newcommand{\RobRTwoCatSciTechMeanDiff}{\ensuremath{-0.013}}
% r2_results.json: by_category.Sports.n
\newcommand{\RobRTwoCatSportsN}{\ensuremath{19{,}162}}
% r2_results.json: by_category.Sports.median_abs_delta
\newcommand{\RobRTwoCatSportsMedianAbsDelta}{\ensuremath{0.020}}
% r2_results.json: by_category.Sports.brier.mean_diff
\newcommand{\RobRTwoCatSportsMeanDiff}{\ensuremath{-0.039}}
% 142123_robust_full_v4/markets.csv: category_refined=Crypto, row count
\newcommand{\CatCryptoNInScope}{\ensuremath{102{,}204}}
% 142123_robust_full_v4/markets.csv: category_refined=Crypto, row count / all rows
\newcommand{\CatCryptoShareInScope}{\ensuremath{14.9}\%}
% 142123_robust_full_v4/markets.csv: category_refined=Crypto, kept=True
\newcommand{\CatCryptoNKept}{\ensuremath{11{,}888}}
% 142123_robust_full_v4/markets.csv: category_refined=Crypto, kept=True / all kept
\newcommand{\CatCryptoShareKept}{\ensuremath{24.8}\%}
% 142123_robust_full_v4/markets.csv: category_refined=Crypto, kept=True / row count
\newcommand{\CatCryptoSurvival}{\ensuremath{11.6}\%}
% 142123_robust_full_v4/markets.csv: category_refined=Crypto, drop_reason=updown
\newcommand{\CatCryptoDropUpdownN}{\ensuremath{60{,}742}}
% 142123_robust_full_v4/markets.csv: category_refined=Crypto, drop_reason=no_outcome
\newcommand{\CatCryptoDropNoOutcomeN}{\ensuremath{7{,}696}}
% 142123_robust_full_v4/markets.csv: category_refined=Crypto, drop_reason=no_price
\newcommand{\CatCryptoDropNoPriceN}{\ensuremath{18{,}433}}
% 142123_robust_full_v4/markets.csv: category_refined=Crypto, drop_reason=stale_price
\newcommand{\CatCryptoDropStalePriceN}{\ensuremath{2{,}721}}
% 142123_robust_full_v4/markets.csv: category_refined=Crypto, drop_reason=below_min_trades
\newcommand{\CatCryptoDropBelowMinTradesN}{\ensuremath{724}}
% 142123_robust_full_v4/markets.csv: category_refined=Culture, row count
\newcommand{\CatCultureNInScope}{\ensuremath{2{,}999}}
% 142123_robust_full_v4/markets.csv: category_refined=Culture, row count / all rows
\newcommand{\CatCultureShareInScope}{\ensuremath{0.4}\%}
% 142123_robust_full_v4/markets.csv: category_refined=Culture, kept=True
\newcommand{\CatCultureNKept}{\ensuremath{1{,}074}}
% 142123_robust_full_v4/markets.csv: category_refined=Culture, kept=True / all kept
\newcommand{\CatCultureShareKept}{\ensuremath{2.2}\%}
% 142123_robust_full_v4/markets.csv: category_refined=Culture, kept=True / row count
\newcommand{\CatCultureSurvival}{\ensuremath{35.8}\%}
% 142123_robust_full_v4/markets.csv: category_refined=Culture, drop_reason=updown
\newcommand{\CatCultureDropUpdownN}{\ensuremath{0}}
% 142123_robust_full_v4/markets.csv: category_refined=Culture, drop_reason=no_outcome
\newcommand{\CatCultureDropNoOutcomeN}{\ensuremath{280}}
% 142123_robust_full_v4/markets.csv: category_refined=Culture, drop_reason=no_price
\newcommand{\CatCultureDropNoPriceN}{\ensuremath{368}}
% 142123_robust_full_v4/markets.csv: category_refined=Culture, drop_reason=stale_price
\newcommand{\CatCultureDropStalePriceN}{\ensuremath{1{,}140}}
% 142123_robust_full_v4/markets.csv: category_refined=Culture, drop_reason=below_min_trades
\newcommand{\CatCultureDropBelowMinTradesN}{\ensuremath{137}}
% 142123_robust_full_v4/markets.csv: category_refined=Finance, row count
\newcommand{\CatFinanceNInScope}{\ensuremath{13{,}370}}
% 142123_robust_full_v4/markets.csv: category_refined=Finance, row count / all rows
\newcommand{\CatFinanceShareInScope}{\ensuremath{1.9}\%}
% 142123_robust_full_v4/markets.csv: category_refined=Finance, kept=True
\newcommand{\CatFinanceNKept}{\ensuremath{5{,}146}}
% 142123_robust_full_v4/markets.csv: category_refined=Finance, kept=True / all kept
\newcommand{\CatFinanceShareKept}{\ensuremath{10.7}\%}
% 142123_robust_full_v4/markets.csv: category_refined=Finance, kept=True / row count
\newcommand{\CatFinanceSurvival}{\ensuremath{38.5}\%}
% 142123_robust_full_v4/markets.csv: category_refined=Finance, drop_reason=updown
\newcommand{\CatFinanceDropUpdownN}{\ensuremath{1{,}574}}
% 142123_robust_full_v4/markets.csv: category_refined=Finance, drop_reason=no_outcome
\newcommand{\CatFinanceDropNoOutcomeN}{\ensuremath{996}}
% 142123_robust_full_v4/markets.csv: category_refined=Finance, drop_reason=no_price
\newcommand{\CatFinanceDropNoPriceN}{\ensuremath{1{,}117}}
% 142123_robust_full_v4/markets.csv: category_refined=Finance, drop_reason=stale_price
\newcommand{\CatFinanceDropStalePriceN}{\ensuremath{2{,}506}}
% 142123_robust_full_v4/markets.csv: category_refined=Finance, drop_reason=below_min_trades
\newcommand{\CatFinanceDropBelowMinTradesN}{\ensuremath{2{,}031}}
% 142123_robust_full_v4/markets.csv: category_refined=Other, row count
\newcommand{\CatOtherNInScope}{\ensuremath{13{,}775}}
% 142123_robust_full_v4/markets.csv: category_refined=Other, row count / all rows
\newcommand{\CatOtherShareInScope}{\ensuremath{2.0}\%}
% 142123_robust_full_v4/markets.csv: category_refined=Other, kept=True
\newcommand{\CatOtherNKept}{\ensuremath{2{,}274}}
% 142123_robust_full_v4/markets.csv: category_refined=Other, kept=True / all kept
\newcommand{\CatOtherShareKept}{\ensuremath{4.7}\%}
% 142123_robust_full_v4/markets.csv: category_refined=Other, kept=True / row count
\newcommand{\CatOtherSurvival}{\ensuremath{16.5}\%}
% 142123_robust_full_v4/markets.csv: category_refined=Other, drop_reason=updown
\newcommand{\CatOtherDropUpdownN}{\ensuremath{1{,}406}}
% 142123_robust_full_v4/markets.csv: category_refined=Other, drop_reason=no_outcome
\newcommand{\CatOtherDropNoOutcomeN}{\ensuremath{1{,}426}}
% 142123_robust_full_v4/markets.csv: category_refined=Other, drop_reason=no_price
\newcommand{\CatOtherDropNoPriceN}{\ensuremath{5{,}493}}
% 142123_robust_full_v4/markets.csv: category_refined=Other, drop_reason=stale_price
\newcommand{\CatOtherDropStalePriceN}{\ensuremath{2{,}044}}
% 142123_robust_full_v4/markets.csv: category_refined=Other, drop_reason=below_min_trades
\newcommand{\CatOtherDropBelowMinTradesN}{\ensuremath{1{,}132}}
% 142123_robust_full_v4/markets.csv: category_refined=Politics, row count
\newcommand{\CatPoliticsNInScope}{\ensuremath{18{,}250}}
% 142123_robust_full_v4/markets.csv: category_refined=Politics, row count / all rows
\newcommand{\CatPoliticsShareInScope}{\ensuremath{2.7}\%}
% 142123_robust_full_v4/markets.csv: category_refined=Politics, kept=True
\newcommand{\CatPoliticsNKept}{\ensuremath{6{,}922}}
% 142123_robust_full_v4/markets.csv: category_refined=Politics, kept=True / all kept
\newcommand{\CatPoliticsShareKept}{\ensuremath{14.4}\%}
% 142123_robust_full_v4/markets.csv: category_refined=Politics, kept=True / row count
\newcommand{\CatPoliticsSurvival}{\ensuremath{37.9}\%}
% 142123_robust_full_v4/markets.csv: category_refined=Politics, drop_reason=updown
\newcommand{\CatPoliticsDropUpdownN}{\ensuremath{56}}
% 142123_robust_full_v4/markets.csv: category_refined=Politics, drop_reason=no_outcome
\newcommand{\CatPoliticsDropNoOutcomeN}{\ensuremath{1{,}527}}
% 142123_robust_full_v4/markets.csv: category_refined=Politics, drop_reason=no_price
\newcommand{\CatPoliticsDropNoPriceN}{\ensuremath{4{,}973}}
% 142123_robust_full_v4/markets.csv: category_refined=Politics, drop_reason=stale_price
\newcommand{\CatPoliticsDropStalePriceN}{\ensuremath{4{,}035}}
% 142123_robust_full_v4/markets.csv: category_refined=Politics, drop_reason=below_min_trades
\newcommand{\CatPoliticsDropBelowMinTradesN}{\ensuremath{737}}
% 142123_robust_full_v4/markets.csv: category_refined=Price Action, row count
\newcommand{\CatPriceActionNInScope}{\ensuremath{192{,}762}}
% 142123_robust_full_v4/markets.csv: category_refined=Price Action, row count / all rows
\newcommand{\CatPriceActionShareInScope}{\ensuremath{28.0}\%}
% 142123_robust_full_v4/markets.csv: category_refined=Price Action, kept=True
\newcommand{\CatPriceActionNKept}{\ensuremath{663}}
% 142123_robust_full_v4/markets.csv: category_refined=Price Action, kept=True / all kept
\newcommand{\CatPriceActionShareKept}{\ensuremath{1.4}\%}
% 142123_robust_full_v4/markets.csv: category_refined=Price Action, kept=True / row count
\newcommand{\CatPriceActionSurvival}{\ensuremath{0.3}\%}
% 142123_robust_full_v4/markets.csv: category_refined=Price Action, drop_reason=updown
\newcommand{\CatPriceActionDropUpdownN}{\ensuremath{191{,}290}}
% 142123_robust_full_v4/markets.csv: category_refined=Price Action, drop_reason=no_outcome
\newcommand{\CatPriceActionDropNoOutcomeN}{\ensuremath{214}}
% 142123_robust_full_v4/markets.csv: category_refined=Price Action, drop_reason=no_price
\newcommand{\CatPriceActionDropNoPriceN}{\ensuremath{373}}
% 142123_robust_full_v4/markets.csv: category_refined=Price Action, drop_reason=stale_price
\newcommand{\CatPriceActionDropStalePriceN}{\ensuremath{183}}
% 142123_robust_full_v4/markets.csv: category_refined=Price Action, drop_reason=below_min_trades
\newcommand{\CatPriceActionDropBelowMinTradesN}{\ensuremath{39}}
% 142123_robust_full_v4/markets.csv: category_refined=Sci-Tech, row count
\newcommand{\CatSciTechNInScope}{\ensuremath{2{,}860}}
% 142123_robust_full_v4/markets.csv: category_refined=Sci-Tech, row count / all rows
\newcommand{\CatSciTechShareInScope}{\ensuremath{0.4}\%}
% 142123_robust_full_v4/markets.csv: category_refined=Sci-Tech, kept=True
\newcommand{\CatSciTechNKept}{\ensuremath{819}}
% 142123_robust_full_v4/markets.csv: category_refined=Sci-Tech, kept=True / all kept
\newcommand{\CatSciTechShareKept}{\ensuremath{1.7}\%}
% 142123_robust_full_v4/markets.csv: category_refined=Sci-Tech, kept=True / row count
\newcommand{\CatSciTechSurvival}{\ensuremath{28.6}\%}
% 142123_robust_full_v4/markets.csv: category_refined=Sci-Tech, drop_reason=updown
\newcommand{\CatSciTechDropUpdownN}{\ensuremath{233}}
% 142123_robust_full_v4/markets.csv: category_refined=Sci-Tech, drop_reason=no_outcome
\newcommand{\CatSciTechDropNoOutcomeN}{\ensuremath{252}}
% 142123_robust_full_v4/markets.csv: category_refined=Sci-Tech, drop_reason=no_price
\newcommand{\CatSciTechDropNoPriceN}{\ensuremath{533}}
% 142123_robust_full_v4/markets.csv: category_refined=Sci-Tech, drop_reason=stale_price
\newcommand{\CatSciTechDropStalePriceN}{\ensuremath{794}}
% 142123_robust_full_v4/markets.csv: category_refined=Sci-Tech, drop_reason=below_min_trades
\newcommand{\CatSciTechDropBelowMinTradesN}{\ensuremath{229}}
% 142123_robust_full_v4/markets.csv: category_refined=Sports, row count
\newcommand{\CatSportsNInScope}{\ensuremath{341{,}911}}
% 142123_robust_full_v4/markets.csv: category_refined=Sports, row count / all rows
\newcommand{\CatSportsShareInScope}{\ensuremath{49.7}\%}
% 142123_robust_full_v4/markets.csv: category_refined=Sports, kept=True
\newcommand{\CatSportsNKept}{\ensuremath{19{,}162}}
% 142123_robust_full_v4/markets.csv: category_refined=Sports, kept=True / all kept
\newcommand{\CatSportsShareKept}{\ensuremath{40.0}\%}
% 142123_robust_full_v4/markets.csv: category_refined=Sports, kept=True / row count
\newcommand{\CatSportsSurvival}{\ensuremath{5.6}\%}
% 142123_robust_full_v4/markets.csv: category_refined=Sports, row count
\newcommand{\SportsNInScope}{\ensuremath{341{,}911}}
% 142123_robust_full_v4/markets.csv: category_refined=Sports, row count / all rows
\newcommand{\SportsShareInScope}{\ensuremath{49.7}\%}
% 142123_robust_full_v4/markets.csv: category_refined=Sports, kept=True / all kept
\newcommand{\SportsShareKept}{\ensuremath{40.0}\%}
% 142123_robust_full_v4/markets.csv: category_refined=Sports, drop_reason=updown
\newcommand{\CatSportsDropUpdownN}{\ensuremath{119}}
% 142123_robust_full_v4/markets.csv: category_refined=Sports, drop_reason=no_outcome
\newcommand{\CatSportsDropNoOutcomeN}{\ensuremath{50{,}327}}
% 142123_robust_full_v4/markets.csv: category_refined=Sports, drop_reason=no_price
\newcommand{\CatSportsDropNoPriceN}{\ensuremath{183{,}999}}
% 142123_robust_full_v4/markets.csv: category_refined=Sports, drop_reason=stale_price
\newcommand{\CatSportsDropStalePriceN}{\ensuremath{66{,}807}}
% 142123_robust_full_v4/markets.csv: category_refined=Sports, drop_reason=below_min_trades
\newcommand{\CatSportsDropBelowMinTradesN}{\ensuremath{21{,}497}}
% m1_results.json: rows.p_1d (main).net_gap_ci_low
\newcommand{\RThreeQOneNetGapCiLow}{\ensuremath{-1.28}\%}
% m1_results.json: rows.p_1d (main).net_gap_ci_high
\newcommand{\RThreeQOneNetGapCiHigh}{\ensuremath{-0.70}\%}
% m1_results.json: rows.p_vwap.row_agreement
\newcommand{\RThreeMOneVwapRuleSame}{yes}
% m1_results.json: rows.p_vwap.n
\newcommand{\RThreeMOneVwapN}{\ensuremath{47{,}948}}
% m1_results.json: rows.p_vwap.chi2
\newcommand{\RThreeMOneVwapChiSq}{\ensuremath{178.0}}
% m1_results.json: rows.p_vwap.bins_off_bh
\newcommand{\RThreeMOneVwapBinsBh}{\ensuremath{6}}
% m1_results.json: rows.p_vwap.net_gap
\newcommand{\RThreeMOneVwapNetGap}{\ensuremath{-1.20}\%}
% m1_results.json: rows.p_vwap.net_gap_ci_low
\newcommand{\RThreeMOneVwapNetGapCiLow}{\ensuremath{-1.50}\%}
% m1_results.json: rows.p_vwap.net_gap_ci_high
\newcommand{\RThreeMOneVwapNetGapCiHigh}{\ensuremath{-0.91}\%}
% m1_results.json: rows.p_mid.n
\newcommand{\RThreeMOneMidN}{\ensuremath{37{,}446}}
% m1_results.json: rows.p_mid.chi2
\newcommand{\RThreeMOneMidChiSq}{\ensuremath{106.9}}
% m1_results.json: rows.p_mid.bins_off_bh
\newcommand{\RThreeMOneMidBinsBh}{\ensuremath{5}}
% m1_results.json: rows.p_mid.net_gap
\newcommand{\RThreeMOneMidNetGap}{\ensuremath{-1.13}\%}
% m1_results.json: rows.p_mid.net_gap_ci_low
\newcommand{\RThreeMOneMidNetGapCiLow}{\ensuremath{-1.48}\%}
% m1_results.json: rows.p_mid.net_gap_ci_high
\newcommand{\RThreeMOneMidNetGapCiHigh}{\ensuremath{-0.78}\%}
% m1_results.json: mean_h
\newcommand{\RThreeMOneMeanHalfSpread}{\ensuremath{1.57}\%}
% m1_results.json: coverage.share_S_both
\newcommand{\RThreeMOneShareBoth}{\ensuremath{78.1}\%}
% m1_results.json: side_line.gap_within_half_spread
\newcommand{\RThreeMOneGapWithinHalfSpread}{yes}
% m2_results.json: yes_no_n
\newcommand{\RThreeMTwoYesNoN}{\ensuremath{30{,}112}}
% m2_results.json: yes_no_chi2
\newcommand{\RThreeMTwoYesNoChiSq}{\ensuremath{100.1}}
% m2_results.json: yes_no_p
\newcommand{\RThreeMTwoYesNoP}{\ensuremath{5.1\times10^{-17}}}
% m2_results.json: yes_no_bins_off_bh
\newcommand{\RThreeMTwoYesNoBinsBh}{\ensuremath{7}}
% m2_results.json: yes_no_net_gap
\newcommand{\RThreeMTwoYesNoNetGap}{\ensuremath{-1.32}\%}
% m2_results.json: yes_no_net_gap_ci_low
\newcommand{\RThreeMTwoYesNoNetGapCiLow}{\ensuremath{-1.61}\%}
% m2_results.json: yes_no_net_gap_ci_high
\newcommand{\RThreeMTwoYesNoNetGapCiHigh}{\ensuremath{-1.03}\%}
% m2_results.json: other_n
\newcommand{\RThreeMTwoOtherN}{\ensuremath{17{,}836}}
% m2_results.json: other_chi2
\newcommand{\RThreeMTwoOtherChiSq}{\ensuremath{15.0}}
% m2_results.json: other_p
\newcommand{\RThreeMTwoOtherP}{\ensuremath{0.133}}
% m2_results.json: other_net_gap
\newcommand{\RThreeMTwoOtherNetGap}{\ensuremath{-0.43}\%}
% m2_results.json: other_net_gap_ci_low
\newcommand{\RThreeMTwoOtherNetGapCiLow}{\ensuremath{-1.05}\%}
% m2_results.json: other_net_gap_ci_high
\newcommand{\RThreeMTwoOtherNetGapCiHigh}{\ensuremath{0.18}\%}
% m2_results.json: yes_no_row_agreement
\newcommand{\RThreeMTwoYesNoRuleSame}{yes}
% m2_results.json: other_row_agreement
\newcommand{\RThreeMTwoOtherRuleSame}{no}
% m3_bootstrap.json: proxies.k1_slug_stem.n_clusters
\newcommand{\RThreeMThreeKOneNClusters}{\ensuremath{16{,}161}}
% m3_bootstrap.json: proxies.k1_slug_stem.net_gap_ci_low
\newcommand{\RThreeMThreeKOneNetGapCiLow}{\ensuremath{-1.41}\%}
% m3_bootstrap.json: proxies.k1_slug_stem.net_gap_ci_high
\newcommand{\RThreeMThreeKOneNetGapCiHigh}{\ensuremath{-0.58}\%}
% m3_bootstrap.json: proxies.k1_slug_stem.x2_cl
\newcommand{\RThreeMThreeKOneChiSq}{\ensuremath{62.8}}
% m3_bootstrap.json: proxies.k1_slug_stem.p
\newcommand{\RThreeMThreeKOneP}{\ensuremath{1.1\times10^{-9}}}
% m3_bootstrap.json: proxies.k2_category_close_hour.n_clusters
\newcommand{\RThreeMThreeKTwoNClusters}{\ensuremath{6{,}181}}
% m3_bootstrap.json: proxies.k2_category_close_hour.net_gap_ci_low
\newcommand{\RThreeMThreeKTwoNetGapCiLow}{\ensuremath{-1.39}\%}
% m3_bootstrap.json: proxies.k2_category_close_hour.net_gap_ci_high
\newcommand{\RThreeMThreeKTwoNetGapCiHigh}{\ensuremath{-0.58}\%}
% m3_bootstrap.json: proxies.k2_category_close_hour.x2_cl
\newcommand{\RThreeMThreeKTwoChiSq}{\ensuremath{58.5}}
% m3_bootstrap.json: proxies.k2_category_close_hour.p
\newcommand{\RThreeMThreeKTwoP}{\ensuremath{7.0\times10^{-9}}}
% m3_bootstrap.json: proxies.k1_slug_stem.row_agreement
\newcommand{\RThreeMThreeKOneRuleSame}{yes}
% m3_bootstrap.json: read_proxy
\newcommand{\RThreeMThreeReadProxy}{k1\_slug\_stem}
% m4_results.json: kept_volume_before_T1d_median
\newcommand{\RThreeMFourKeptMedianVolumeBeforeDay}{\ensuremath{1{,}790.0}}
% m4_results.json: kept_duration_days_median
\newcommand{\RThreeMFourKeptMedianDuration}{\ensuremath{7.0}}
% m4_results.json: kept_n_trades_median
\newcommand{\RThreeMFourKeptMedianTrades}{\ensuremath{292.5}}
% m4_results.json: eligible_dropped_pooled_volume_before_T1d_median
\newcommand{\RThreeMFourDroppedMedianVolumeBeforeDay}{\ensuremath{69.0}}
% m4_results.json: eligible_dropped_pooled_duration_days_median
\newcommand{\RThreeMFourDroppedMedianDuration}{\ensuremath{6.4}}
% m4_results.json: eligible_dropped_pooled_n_trades_median
\newcommand{\RThreeMFourDroppedMedianTrades}{\ensuremath{17}}
% m4_results.json: stale_q1_net_gap
\newcommand{\RThreeMFourStaleNetGap}{\ensuremath{-0.35}\%}
% m4_results.json: stale_q1_net_gap_ci_low
\newcommand{\RThreeMFourStaleNetGapCiLow}{\ensuremath{-0.48}\%}
% m4_results.json: stale_q1_net_gap_ci_high
\newcommand{\RThreeMFourStaleNetGapCiHigh}{\ensuremath{-0.22}\%}
% m5_calibration_fit.json: a
\newcommand{\RThreeMFiveA}{\ensuremath{-0.089}}
% m5_calibration_fit.json: a_wald_ci[0]
\newcommand{\RThreeMFiveAWaldLow}{\ensuremath{-0.116}}
% m5_calibration_fit.json: a_wald_ci[1]
\newcommand{\RThreeMFiveAWaldHigh}{\ensuremath{-0.061}}
% m5_calibration_fit.json: b
\newcommand{\RThreeMFiveB}{\ensuremath{1.028}}
% m5_calibration_fit.json: b_wald_ci[0]
\newcommand{\RThreeMFiveBWaldLow}{\ensuremath{1.005}}
% m5_calibration_fit.json: b_wald_ci[1]
\newcommand{\RThreeMFiveBWaldHigh}{\ensuremath{1.050}}
% m5_calibration_fit.json: brier
\newcommand{\RThreeMFiveBrier}{\ensuremath{0.107}}
% m5_calibration_fit.json: climatology
\newcommand{\RThreeMFiveClimatologyBrier}{\ensuremath{0.232}}
% m5_calibration_fit.json: brier_skill_score
\newcommand{\RThreeMFiveSkill}{\ensuremath{0.538}}
% m5_calibration_fit.json: majority_accuracy
\newcommand{\RThreeMFiveMajorityAccuracy}{\ensuremath{63.4}\%}
% m5_calibration_fit.json: accuracy_at_0.5
\newcommand{\RThreeMFiveAccuracyHalf}{\ensuremath{83.4}\%}
% m5_calibration_fit.json: delta_table[].delta, largest absolute value
\newcommand{\RThreeMFiveMaxAbsDelta}{\ensuremath{2.22}\%}
% m5_calibration_fit.json: delta_table[].practically_sized, count of True
\newcommand{\RThreeMFivePracticalPrices}{\ensuremath{0}}
% m5_calibration_fit.json: delta_table, length
\newcommand{\RThreeMFiveGridPrices}{\ensuremath{7}}
% minor4_dedup.json: duplicates.share_dup_all_columns
\newcommand{\RThreeMinorFourDupShare}{\ensuremath{3.1}\%}
% minor4_dedup.json: kept_set.n_removed
\newcommand{\RThreeMinorFourRemoved}{\ensuremath{139}}
% minor4_dedup.json: q1.chi2
\newcommand{\RThreeMinorFourChiSq}{\ensuremath{67.4}}
% minor4_dedup.json: q1.net_gap
\newcommand{\RThreeMinorFourNetGap}{\ensuremath{-0.99}\%}
% minor4_dedup.json: kept_set.n_out
\newcommand{\RThreeMinorFourKeptN}{\ensuremath{47{,}809}}
% minor4_dedup.json: q1.net_gap_ci_low
\newcommand{\RThreeMinorFourNetGapCiLow}{\ensuremath{-1.28}\%}
% minor4_dedup.json: q1.net_gap_ci_high
\newcommand{\RThreeMinorFourNetGapCiHigh}{\ensuremath{-0.69}\%}
% minor4_dedup.json: q1.row_agreement
\newcommand{\RThreeMinorFourRuleSame}{yes}
% minor4_dedup.json: q2.maker.top1_share.delta_pp (percentage points)
\newcommand{\RThreeMinorFourMakerTopOneDeltaPp}{\ensuremath{-0.16}}
% minor4_dedup.json: q2.maker.top20_share.delta_pp (percentage points)
\newcommand{\RThreeMinorFourMakerTopTwentyDeltaPp}{\ensuremath{-0.02}}
% minor4_dedup.json: q2.taker.top1_share.delta_pp (percentage points)
\newcommand{\RThreeMinorFourTakerTopOneDeltaPp}{\ensuremath{-0.26}}
% minor4_dedup.json: q2.taker.top20_share.delta_pp (percentage points)
\newcommand{\RThreeMinorFourTakerTopTwentyDeltaPp}{\ensuremath{-0.03}}
% peer_a_results.json: cut_low_mid
\newcommand{\RThreePeerACutLowMid}{\ensuremath{498.5}}
% peer_a_results.json: cut_mid_high
\newcommand{\RThreePeerACutMidHigh}{\ensuremath{6{,}207.4}}
% peer_a_results.json: tercile_low_n
\newcommand{\RThreePeerALowN}{\ensuremath{15{,}983}}
% peer_a_results.json: tercile_low_chi2
\newcommand{\RThreePeerALowChiSq}{\ensuremath{40.3}}
% peer_a_results.json: tercile_low_p_bh
\newcommand{\RThreePeerALowBhP}{\ensuremath{2.2\times10^{-5}}}
% peer_a_results.json: tercile_low_net_gap
\newcommand{\RThreePeerALowNetGap}{\ensuremath{-1.12}\%}
% peer_a_results.json: tercile_low_net_gap_ci_low
\newcommand{\RThreePeerALowNetGapCiLow}{\ensuremath{-1.66}\%}
% peer_a_results.json: tercile_low_net_gap_ci_high
\newcommand{\RThreePeerALowNetGapCiHigh}{\ensuremath{-0.59}\%}
% peer_a_results.json: tercile_low_row_agreement
\newcommand{\RThreePeerALowRuleSame}{yes}
% peer_a_results.json: tercile_mid_n
\newcommand{\RThreePeerAMidN}{\ensuremath{15{,}982}}
% peer_a_results.json: tercile_mid_chi2
\newcommand{\RThreePeerAMidChiSq}{\ensuremath{45.0}}
% peer_a_results.json: tercile_mid_p_bh
\newcommand{\RThreePeerAMidBhP}{\ensuremath{6.4\times10^{-6}}}
% peer_a_results.json: tercile_mid_net_gap
\newcommand{\RThreePeerAMidNetGap}{\ensuremath{-1.43}\%}
% peer_a_results.json: tercile_mid_net_gap_ci_low
\newcommand{\RThreePeerAMidNetGapCiLow}{\ensuremath{-1.97}\%}
% peer_a_results.json: tercile_mid_net_gap_ci_high
\newcommand{\RThreePeerAMidNetGapCiHigh}{\ensuremath{-0.89}\%}
% peer_a_results.json: tercile_mid_row_agreement
\newcommand{\RThreePeerAMidRuleSame}{yes}
% peer_a_results.json: tercile_high_n
\newcommand{\RThreePeerAHighN}{\ensuremath{15{,}983}}
% peer_a_results.json: tercile_high_chi2
\newcommand{\RThreePeerAHighChiSq}{\ensuremath{14.2}}
% peer_a_results.json: tercile_high_p_bh
\newcommand{\RThreePeerAHighBhP}{\ensuremath{0.166}}
% peer_a_results.json: tercile_high_net_gap
\newcommand{\RThreePeerAHighNetGap}{\ensuremath{-0.42}\%}
% peer_a_results.json: tercile_high_net_gap_ci_low
\newcommand{\RThreePeerAHighNetGapCiLow}{\ensuremath{-0.86}\%}
% peer_a_results.json: tercile_high_net_gap_ci_high
\newcommand{\RThreePeerAHighNetGapCiHigh}{\ensuremath{0.03}\%}
% peer_a_results.json: tercile_high_row_agreement
\newcommand{\RThreePeerAHighRuleSame}{no}
% peer_a_results.json: high_vs_low.diff
\newcommand{\RThreePeerADiff}{\ensuremath{0.71}\%}
% peer_a_results.json: high_vs_low.ci_low
\newcommand{\RThreePeerADiffCiLow}{\ensuremath{0.01}\%}
% peer_a_results.json: high_vs_low.ci_high
\newcommand{\RThreePeerADiffCiHigh}{\ensuremath{1.40}\%}
% peer_a_results.json: high_vs_low.p
\newcommand{\RThreePeerADiffP}{\ensuremath{0.046}}
% peer_b_results.json: year_to_period values, period 1 in file order
\newcommand{\RThreePeerBPeriodOneLabel}{2022-2024}
% peer_b_results.json: period_2022-2024_n
\newcommand{\RThreePeerBPeriodOneN}{\ensuremath{3{,}444}}
% peer_b_results.json: period_2022-2024_chi2
\newcommand{\RThreePeerBPeriodOneChiSq}{\ensuremath{25.1}}
% peer_b_results.json: period_2022-2024_p_bh
\newcommand{\RThreePeerBPeriodOneBhP}{\ensuremath{0.008}}
% peer_b_results.json: period_2022-2024_net_gap
\newcommand{\RThreePeerBPeriodOneNetGap}{\ensuremath{-1.71}\%}
% peer_b_results.json: period_2022-2024_net_gap_ci_low
\newcommand{\RThreePeerBPeriodOneNetGapCiLow}{\ensuremath{-2.81}\%}
% peer_b_results.json: period_2022-2024_net_gap_ci_high
\newcommand{\RThreePeerBPeriodOneNetGapCiHigh}{\ensuremath{-0.62}\%}
% peer_b_results.json: period_2022-2024_row_agreement
\newcommand{\RThreePeerBPeriodOneRuleSame}{yes}
% peer_b_results.json: year_to_period values, period 2 in file order
\newcommand{\RThreePeerBPeriodTwoLabel}{2025}
% peer_b_results.json: period_2025_n
\newcommand{\RThreePeerBPeriodTwoN}{\ensuremath{17{,}816}}
% peer_b_results.json: period_2025_chi2
\newcommand{\RThreePeerBPeriodTwoChiSq}{\ensuremath{76.3}}
% peer_b_results.json: period_2025_p_bh
\newcommand{\RThreePeerBPeriodTwoBhP}{\ensuremath{7.8\times10^{-12}}}
% peer_b_results.json: period_2025_net_gap
\newcommand{\RThreePeerBPeriodTwoNetGap}{\ensuremath{-1.36}\%}
% peer_b_results.json: period_2025_net_gap_ci_low
\newcommand{\RThreePeerBPeriodTwoNetGapCiLow}{\ensuremath{-1.84}\%}
% peer_b_results.json: period_2025_net_gap_ci_high
\newcommand{\RThreePeerBPeriodTwoNetGapCiHigh}{\ensuremath{-0.89}\%}
% peer_b_results.json: period_2025_row_agreement
\newcommand{\RThreePeerBPeriodTwoRuleSame}{yes}
% peer_b_results.json: year_to_period values, period 3 in file order
\newcommand{\RThreePeerBPeriodThreeLabel}{2026}
% peer_b_results.json: period_2026_n
\newcommand{\RThreePeerBPeriodThreeN}{\ensuremath{26{,}688}}
% peer_b_results.json: period_2026_chi2
\newcommand{\RThreePeerBPeriodThreeChiSq}{\ensuremath{15.8}}
% peer_b_results.json: period_2026_p_bh
\newcommand{\RThreePeerBPeriodThreeBhP}{\ensuremath{0.105}}
% peer_b_results.json: period_2026_net_gap
\newcommand{\RThreePeerBPeriodThreeNetGap}{\ensuremath{-0.65}\%}
% peer_b_results.json: period_2026_net_gap_ci_low
\newcommand{\RThreePeerBPeriodThreeNetGapCiLow}{\ensuremath{-1.04}\%}
% peer_b_results.json: period_2026_net_gap_ci_high
\newcommand{\RThreePeerBPeriodThreeNetGapCiHigh}{\ensuremath{-0.25}\%}
% peer_b_results.json: period_2026_row_agreement
\newcommand{\RThreePeerBPeriodThreeRuleSame}{no}
% wallet PnL run directory name
\newcommand{\PnlRun}{156432\_pnl\_wallets\_v1}
% pnl_summary.json (run 156432), scopes.kept: n_wallets
\newcommand{\PnlNWallets}{\ensuremath{1{,}569{,}986}}
% pnl_summary.json (run 156432), scopes.kept: n_winners (pnl_net > 1e-6)
\newcommand{\PnlNWinners}{\ensuremath{663{,}072}}
% pnl_summary.json (run 156432), scopes.kept: n_losers (pnl_net < -1e-6)
\newcommand{\PnlNLosers}{\ensuremath{871{,}667}}
% pnl_gross_summary.json (run 156432), kept: won_M, USDC millions, before fees (pnl_gross > 0)
\newcommand{\PnlGrossWonM}{\ensuremath{389.8}}
% pnl_gross_summary.json (run 156432), kept: lost_M, USDC millions, before fees (pnl_gross < 0)
\newcommand{\PnlGrossLostM}{\ensuremath{389.8}}
% pnl_summary.json (run 156432), scopes.kept: median_pnl_net, USD per wallet
\newcommand{\PnlMedianNet}{\ensuremath{-0.20}}
% pnl_summary.json (run 156432), scopes.kept: top_pct_of_all_wallets_share_of_won.1%
\newcommand{\PnlTopOnePctShareWon}{\ensuremath{88.0}\%}
% pnl_summary.json (run 156432), scopes.kept: quantiles_pnl_net.1%, USD
\newcommand{\PnlQuantileOne}{\ensuremath{-3{,}145.02}}
% pnl_summary.json (run 156432), scopes.kept: quantiles_pnl_net.99%, USD
\newcommand{\PnlQuantileNinetyNine}{\ensuremath{1{,}968.24}}
% pnl_summary.json (run 156432), checks: n_identity_failures
\newcommand{\PnlNIdentityFailures}{\ensuremath{0}}
% pnl_summary.json (run 156432), checks: n_markets_checked
\newcommand{\PnlNMarketsChecked}{\ensuremath{548{,}340}}
% pnl_summary.json (run 156432), checks: n_kept_not_usable
\newcommand{\PnlNKeptNotUsable}{\ensuremath{1{,}014}}
% pnl_summary.json (run 156432), checks: share_short_positions_kept
\newcommand{\PnlShareShortPositions}{\ensuremath{59.0}\%}
% p_T calibration run directory name
\newcommand{\RTwoCloseRun}{local\_r2close\_v14}
% results.json (run local_r2close_v14): p_1d.chi_square
\newcommand{\RTwoCloseOneDayChiSq}{\ensuremath{68.0}}
% results.json (run local_r2close_v14): p_1d.p
\newcommand{\RTwoCloseOneDayP}{\ensuremath{1.1\times10^{-10}}}
% results.json (run local_r2close_v14): p_1d.bins_bh
\newcommand{\RTwoCloseOneDayBinsBh}{\ensuremath{5}}
% results.json (run local_r2close_v14): p_1d.bins_bonferroni
\newcommand{\RTwoCloseOneDayBinsBonferroni}{\ensuremath{4}}
% results.json (run local_r2close_v14): p_1d.bins_negative
\newcommand{\RTwoCloseOneDayBinsNeg}{\ensuremath{9}}
% results.json (run local_r2close_v14): p_1d.net_gap_pp
\newcommand{\RTwoCloseOneDayNetGap}{\ensuremath{-0.99}\%}
% bins_p_1d.csv: bin 0 n
\newcommand{\RTwoCloseOneDayBinZeroN}{\ensuremath{18{,}209}}
% results.json (run local_r2close_v14): p_T.chi_square
\newcommand{\RTwoCloseCloseChiSq}{\ensuremath{53.6}}
% results.json (run local_r2close_v14): p_T.p
\newcommand{\RTwoCloseCloseP}{\ensuremath{5.7\times10^{-8}}}
% results.json (run local_r2close_v14): p_T.bins_bh
\newcommand{\RTwoCloseCloseBinsBh}{\ensuremath{5}}
% results.json (run local_r2close_v14): p_T.bins_bonferroni
\newcommand{\RTwoCloseCloseBinsBonferroni}{\ensuremath{3}}
% results.json (run local_r2close_v14): p_T.bins_negative
\newcommand{\RTwoCloseCloseBinsNeg}{\ensuremath{9}}
% results.json (run local_r2close_v14): p_T.net_gap_pp
\newcommand{\RTwoCloseCloseNetGap}{\ensuremath{-0.65}\%}
% bins_p_T.csv: bin 0 n
\newcommand{\RTwoCloseCloseBinZeroN}{\ensuremath{21{,}897}}


% --- Design constants and literature facts. Not measured by our runs. ---
% src: S-PAPER-3 (design constants)
\newcommand{\NBins}{10}
\newcommand{\CiLevel}{95\%}
\newcommand{\DaysShort}{1}
\newcommand{\MaxStaleHours}{24}
\newcommand{\TopOneLabel}{1\%}
\newcommand{\TopFiveLabel}{5\%}
\newcommand{\TopTenLabel}{10\%}
\newcommand{\TopTwentyLabel}{20\%}
\newcommand{\HookPct}{70\%}
\newcommand{\pT}{\ensuremath{p_{T}}}
\newcommand{\DeltaA}{0.10}
\newcommand{\DeltaB}{0.25}
\newcommand{\DeltaC}{0.50}
\newcommand{\PayoutDollar}{\$1}
\newcommand{\HookPrice}{\$0.70}
\newcommand{\FavoriteLine}{0.5}
\newcommand{\pOne}{\ensuremath{p_{\DaysShort\mathrm{d}}}}
% src: S-PAPER-1, S-PAPER-2 (facts about the dataset paper)
\newcommand{\PaperTrades}{1.2 billion}
\newcommand{\PaperMarkets}{1.3 million}
\newcommand{\PaperVolume}{\$61 billion}
\newcommand{\PaperPeriod}{November 2022 to April 2026}
\newcommand{\PaperBrierMarkets}{1{,}019}
\newcommand{\PaperTopShare}{84.1\%}
\newcommand{\PaperCheapGap}{0.2 to 0.3 cents}
\newcommand{\PaperCheapPrice}{30 cents}
% Run folder names are identifiers, not measured numbers.
\newcommand{\CcdfRun}{145892\_ccdf\_v6}
\newcommand{\WalletDistRun}{145893\_walletdist\_v7}
\newcommand{\QqRun}{145897\_q3figs\_v8}
\newcommand{\RThreeRun}{150310\_r3\_review\_fixes\_v6}
\newcommand{\RThreeFigRun}{local\_r3figs\_v12}
\newcommand{\RThreeCodeVersion}{v6}
\newcommand{\RThreePreregCommit}{aef0808}

% --- Table row helpers: pull one row of macros by name ---
\newcommand{\binrow}[2]{#1 & \csname Bin#2N\endcsname & \csname Bin#2Price\endcsname
  & \csname Bin#2Obs\endcsname & [\csname Bin#2ZLow\endcsname, \csname Bin#2ZHigh\endcsname]
  & \csname Bin#2Diff\endcsname & \csname Bin#2BhP\endcsname}
\title{When Polymarket Says \HookPct, Does It Happen \HookPct{} of the Time?\\[2pt]
\large Prediction Accuracy, Trading Concentration, and Statistical Distributions in Prediction Markets}
\author{Group 1 $|$ Team Brier: Taek Lim, Zhugang Liu (Tony), Francisco Cortez\\
CS 6344 Introduction to Data Science}
\date{}

\begin{document}
\maketitle

\begin{abstract}
% src: S-PAPER-1, S-PAPER-3, S-PAPER-5, S-PAPER-10
A Polymarket YES share pays one dollar if the event happens.
So a price of \HookPct{} is read as a \HookPct{} chance.
We test this on resolved Yes/No markets from the Polymarket-v1 database.
We keep \NMarketsKept{} of \NMarketsInScope{} markets in scope.
We read each market's price \DaysShort{} day before its scheduled close.
The pre-registered chi-square test says the prices are not calibrated as a whole
($\chi^2=\QOneChiSq$, $p=\QOneP$).
The observed rate of outcome 1, the first-listed outcome, lies below the mean price in \GapNegBinsMain{} of \GapBinsTotal{} bins.
The net gap is \GapNetDiffMain.
The largest per-bin gap is \LargestGapAbs.
In the kept markets, the top \TopTwentyLabel{} of wallets hold \MakerTopTwentyShare{} of maker USDC volume,
far above the 80/20 reference of \TopTwentyReference.
Mean volume is far above median volume.
Trade sizes are strongly right skewed.
The Q1 effect is small and runs one way.
We set no size line before the results.
We describe Q2 and Q3 without a verdict.
A payout check leaves the Q1 verdict unchanged.
\end{abstract}

\section{Introduction}\label{sec:intro}
% src: S-PAPER-1
A Polymarket YES share pays \PayoutDollar{} if the event happens.
At a price of \HookPrice, the price is read as a \HookPct{} chance.
Is that chance right?
If we collect every event priced at \HookPct, about \HookPct{} of them should have happened.
The Polymarket-v1 database records every trade on Polymarket's first exchange.
It holds \PaperTrades{} trades and \PaperMarkets{} markets from \PaperPeriod{} \citep{qin2026polymarketv1}.
So we can check each price against what happened.

This paper reports the first stage of a staged study.
% src: S-PAPER-8
It answers three questions.
The first question is Q1: is the crowd right?
% src: S-PAPER-1
We test calibration one price range at a time.
We use binomial and chi-square tests, confidence intervals, and the mean squared residual of the price.
We also ask whether cheap longshot bets are a trap.
The second question is Q2: who is the crowd?
We ask whether a small group of wallets does most of the trading.
We use top shares of wallet volume, the 80/20 rule, and a power-law check.
We also group wallets by k-means, as a side analysis.
The third question is Q3: what does betting look like?
We describe trade sizes, wallet volumes, and prices with histograms and distribution fits.

% src: S-PAPER-3, S-PAPER-5
Our contributions are these.
We fix the evaluation rules before we see any Q1 result.
We report one rule change made afterwards.
We score all \NMarketsKept{} kept markets, not a small panel.
We report an effect size next to every p-value, such as the per-bin gap.
The text uses the word wallet for an address on the exchange.

\section{Related Work}\label{sec:related}
% src: S-PAPER-1
We build on two references from our proposal.
Snowberg and Wolfers study horse racing \citep{snowberg2010favorite}.
They show that longshots win less often than their odds say.
We test the same gap on Polymarket.

The dataset paper reports Polymarket results on cheap tokens \citep[Table~3]{qin2026polymarketv1}.
Tokens priced under \PaperCheapPrice{} pay out \PaperCheapGap{} less than their price.
This is counted by trade.
We count markets, not trades.
We read this result as the same direction as the horse-racing bias.

The dataset paper also scores \PaperBrierMarkets{} markets with a squared-error score \citep{qin2026polymarketv1}.
We score all resolved Yes/No markets that pass our filters.
Finally, the paper reports that the top \TopOneLabel{} of maker addresses hold \PaperTopShare{}
of maker volume \citep[Table~5]{qin2026polymarketv1}.
We ask whether a few wallets dominate volume in our kept markets.
% src: S-PAPER-11
Five further references place our tests in context.
Wolfers and Zitzewitz review prediction markets \citep{wolfers2004prediction}.
They report that market forecasts are typically fairly accurate and usually beat moderately sophisticated benchmarks.
We test this accuracy as calibration of Polymarket prices.
Page and Clemen find that prediction markets are reasonably well calibrated when time to expiration is short \citep{page2013prediction}.
They find prices significantly biased for events farther in the future, with a favourite/longshot bias.
Our longshot check looks at the same bias.
Le uses 353 million trades on 429{,}000 binary contracts from Kalshi and Polymarket \citep{le2026decomposing}.
The preprint finds that calibration varies by event domain, time to resolution, and trade size.
We test all kept markets together and do not split by domain or trade size.

Gneiting and Raftery name sharpness subject to calibration as the goal of probabilistic forecasting \citep{gneiting2007strictly}.
They list the Brier score among the proper scoring rules.
Our mean squared residual of the price is this score.
Clauset, Shalizi, and Newman fit power laws by maximum likelihood \citep{clauset2009power}.
They test the fit with the Kolmogorov-Smirnov statistic and with likelihood-ratio tests against alternatives.
On 24 real datasets, the power law fits some and is ruled out for others.
Like them, we measure the fit with the KS statistic.
We fix the lower bound at the top \TopTenLabel{} of values and run no likelihood-ratio test.
\todo{owner: confirm the added references and add any others}

\section{Data}\label{sec:data}
\subsection{Source and scope}\label{sec:data-source}
% src: S-PAPER-1, S-PAPER-2
We use the \texttt{daily\_aligned/} layer of Polymarket-v1.
This layer holds Standard Binary markets (\texttt{neg\_risk} is false).
Its relayer records are already removed \citep{polymarketv1data}.
We keep markets that have a winning outcome label.
We drop Up/Down markets by the default rule in our code.
The dataset paper says they settle by a fundamentally different mechanism and should not be pooled with other markets \citep{qin2026polymarketv1}.
We take the price from the \texttt{p\_event} column, which is on the YES scale.
The data are at \url{https://huggingface.co/datasets/TimeSeventeen/Polymarket-v1} \citep{polymarketv1data}.
The dataset card states the license CC-BY-4.0.
The paper that describes the data is \citep{qin2026polymarketv1}.

\subsection{Filters and drop reasons}\label{sec:data-filters}
% src: S-PAPER-5
In the full run, \NMarketsInScope{} markets are in scope.
We keep \NMarketsKept{} of them (\PctMarketsKept).
\Cref{fig:drops} counts each dropped market once, by the first rule it fails.
The rules apply in the order of the bars.
Up/Down markets are \NUpdown{} (\PctUpdown).
\NNoOutcome{} markets have no usable winner label (\PctNoOutcome).
\NNoPrice{} have no price at or before $T-\DaysShort$ day (\PctNoPrice).
\NStalePrice{} have a price older than \MaxStaleHours{} hours (\PctStalePrice).
\NBelowMinTrades{} have fewer than \MinTrades{} trades before $T-\DaysShort$ day (\PctBelowMinTrades).

\begin{figure}[H]
\centering
\figorbox{figures/drop_reasons}{\linewidth}
\caption{Markets by first failing rule. Why: it shows which rule removes the most markets. Source: run \RunId.}
\label{fig:drops}
\end{figure}

\subsection{Market and wallet tables}\label{sec:data-tables}
% src: S-PAPER-2, S-PAPER-5
We build one market table and one wallet table.
The market table has one row per kept market.
It holds the category, trade count, volume in USDC, the price \pOne, and the outcome $y\in\{0,1\}$.
The wallet table has one row per wallet and role (maker or taker).
It holds volume in USDC, trade count, market count, and mean trade price.
A wallet is one address.
The wallet table has \NWalletRoleRows{} wallet-role rows from the kept markets.
The kept markets hold \NTradesKeptMarkets{} trades.
We keep duplicate rows because no row identifier separates repeat fills from copies.

\section{Methods}\label{sec:methods}
\subsection{Evaluation price and pre-registration}\label{sec:methods-prereg}
% src: S-PAPER-2, S-PAPER-3
Each market has a scheduled close time $T$.
We set $T$ to the \texttt{close\_at} field.
Trades can continue after the close, often once the answer is known.
A price taken after the close can then contain the answer.
Then calibration looks better than it is.
In a one-month pilot of \PilotNMarkets{} markets, \PilotShareMarketsAfterClose{} had a trade after \texttt{close\_at}.
Those late trades are \PilotShareVolumeAfterClose{} of the pilot volume.
So we use the scheduled close and not the last trade.

\Cref{fig:timeline} shows where the evaluation price sits.
The evaluation price \pOne{} is a mean price.
It averages the trades at the last timestamp at or before $T-\DaysShort$ day.
We drop a market if this price is older than \MaxStaleHours{} hours.
We also drop a market with fewer than \MinTrades{} trades before that time.
A price from a nearly empty market is not a crowd view.
We report Q1 with and without this filter.

\begin{figure}[H]
\centering
\begin{tikzpicture}[font=\small]
\draw[thick,-{Stealth}] (0,0) -- (12.2,0);
\draw[thick] (0.2,-0.12) -- (0.2,0.12) node[above] {market opens};
\draw[thick,blue!60!black] (7.0,-0.2) -- (7.0,0.2) node[above] {$T-\DaysShort$ day};
\draw[thick] (9.2,-0.12) -- (9.2,0.12) node[above] {$T$ = \texttt{close\_at}};
\foreach \x in {10.0,10.5,11.0,11.5} \draw[red!70!black,thick] (\x,-0.15) -- (\x,0.15);
\node[below=6pt,blue!60!black,align=center] at (7.0,0) {evaluation price \pOne};
\node[below=6pt,red!70!black,align=center] at (10.75,0) {trades after the close};
\end{tikzpicture}
\caption{Where the evaluation price sits in a market's life (illustrative). Why: a price taken after the close can already contain the answer. Source: pre-registration, \texttt{analysis/prereg\_mvp1.md}.}
\label{fig:timeline}
\end{figure}

The criteria were fixed before any Q1 result.
\Cref{tab:prereg} lists them.
The verdict applies to the full-data run with the same parameters.

\begin{table}[H]
\centering
\caption{Pre-registered choices and Q1 criteria. Why: fixing the rules first keeps the result from choosing the rule. Source: \texttt{analysis/prereg\_mvp1.md}.}
\label{tab:prereg}
\footnotesize
\begin{tabularx}{\linewidth}{@{}lX@{}}
\toprule
Item & Value \\
\midrule
Close time $T$ & \texttt{close\_at} \\
\midrule
Evaluation price & mean price at the last timestamp (after the amendment below) at or before $T-\DaysShort$ day, at most \MaxStaleHours{} hours old \\
\midrule
Minimum trades & \MinTrades{} before $T-\DaysShort$ day. Q1 is also reported without it \\
\midrule
Bins and level & \NBins{} equal-width bins, $\alpha=\Alpha$ \\
\midrule
Primary metric & chi-square goodness of fit over the \NBins{} bins. The line is $p<\Alpha$ \\
\midrule
Per-bin call & binomial test per bin, Benjamini-Hochberg (BH) corrected at \Alpha{} \\
\midrule
Effect size & observed rate minus mean price per bin, with a \CiLevel{} z interval here. Also the mean squared residual of the price \\
\midrule
Failure condition & no bin off after BH and chi-square $p\ge\Alpha$: no evidence of miscalibration at this $n$ \\
\bottomrule
\end{tabularx}
\end{table}

% src: S-PAPER-3
One rule changed after we saw a full-data Q1 result.
The original rule was the last price at or before $T-\DaysShort$ day.
Many trades share one block timestamp, so this rule left their order undefined.
The new rule is the mean price of the trades at the last timestamp.
We made the change for determinism, not because of the result.
It came after we saw a full-data Q1 result, so we report it as an amendment.
The amendment in the pre-registration file records that the verdict did not change.
All numbers in this paper come from the runs with the amended rule.
This paper uses an effect size and an interval that differ from the pre-registration: per-bin gaps and the z interval.
It drops the Gini coefficient named in our proposal.
Both changes keep the analysis within the lecture contents.
The verdict rule and the verdict sentence are unchanged.
These methods are outside the lectures: the Lorenz curve, Wilson intervals, the cluster bootstrap chi-square, logistic calibration, and the Brier skill score.
All other methods come from the CS 6344 course lectures.

\subsection{Q1: calibration}\label{sec:methods-q1}
% src: S-PAPER-2, S-PAPER-3
We treat each outcome as a Bernoulli trial with success probability equal to the price.
Outcome 1 is the first-listed outcome of a market, labeled Yes in \OutcomeOneYesShare{} of kept markets.
So the outcome-1 count in a bin is approximately binomial.
We form \NBins{} equal-width price bins, and we use equal-frequency bins as a check.
In bin $b$, $n_b$ is the market count, $\bar p_b$ the mean price, and $O_b$ the outcome-1 count.
The null hypothesis is that the outcome-1 rate equals the mean price.
The effect size is the per-bin gap, the observed rate minus the mean price.
The expected outcome-1 count is $E_b=n_b\bar p_b$.
We test each bin with a binomial test.
We give a \CiLevel{} z interval for the observed rate $\hat r_b=O_b/n_b$, which is $\hat r_b\pm1.96\sqrt{\hat r_b(1-\hat r_b)/n_b}$.
We correct the \NBins{} tests with Bonferroni and with the Benjamini-Hochberg (BH) procedure \citep{benjamini1995controlling}.
The overall test is the goodness-of-fit form of the chi-square test: observed against expected counts, with the expected counts taken from the prices.
It runs on the bin by \{outcome 1, outcome 0\} table:
\begin{equation}
\chi^2=\sum_{b}\left[\frac{(O_b-E_b)^2}{E_b}+\frac{(O_b-E_b)^2}{n_b-E_b}\right].
\label{eq:chisq}
\end{equation}
We also report the mean squared residual of the price, $\frac1n\sum(p-y)^2$, where $n$ is the total market count.
The pre-registration calls it the Brier score \citep{brier1950verification}.
We add the log loss and the mean entropy $H(p)$ in bits.
We treat the market as a classifier that says outcome 1 wins when $p>\FavoriteLine$.
We compare the price distributions of markets where outcome 1 won and markets where outcome 0 won with a Kolmogorov-Smirnov (KS) test.
The KS statistic $D$ is the effect size.
For the longshot check we look at the two cheapest bins.

\subsection{Q2: wallet concentration}\label{sec:methods-q2}
% src: S-PAPER-2, S-PAPER-10
We work with wallet volume in USDC, by role.
Volume is the USDC sum over every trade in the kept markets.
It includes trades after $T-\DaysShort$ day.
We did not measure trade-count shares.
We report the share of role volume held by the top \TopOneLabel, \TopFiveLabel, \TopTenLabel, and \TopTwentyLabel{} of wallets.
We compare the top \TopTwentyLabel{} share with the 80/20 reference.
This is descriptive, with no pre-registered verdict line.
We draw the complementary CDF (CCDF) of wallet volume on log-log axes.
We sample its ranks from the top, so the largest wallets get points.
We also draw the wallet distribution, the volume distribution, and the Lorenz curve of each role.
Each is split at the top \TopTwentyLabel{} of wallets.
The Lorenz curve is outside the lectures.
We use it only to show how the top shares build up.
As a pipeline check, we compare our top \TopOneLabel{} maker share with the dataset paper's.
We test the log volume of makers against takers with a Welch t-test and a KS test.
The effect size is the difference of the mean log10 volume.
We report Spearman's $\rho$ between trade count and volume.

% src: S-PAPER-2, S-PAPER-6
As a side analysis, we run k-means on wallet features.
We take the logarithm of the heavy-tailed features, then z-score them.
We choose $K$ by a fixed rule on the elbow curve.
The rule takes the point farthest from the line between the first and last points.
We test cluster against topic category with the chi-square test of independence .
Its effect size is the largest observed minus expected share in one cell.
This analysis is exploratory.

\subsection{Q3: distributions}\label{sec:methods-q3}
% src: S-PAPER-2, S-PAPER-5
We summarise price, trade size, and wallet volume with their mean and median.
We plot histograms of the price \pOne{} and of trade size.
We draw QQ plots of trade size and of wallet volume, before and after a log transform.
We fit a normal, an exponential, and a power law to trade size and wallet volume.
The power law is fitted on the top \TopTenLabel{} of values only, above a lower bound $x_{\min}$.
We measure each fit with the KS statistic $D$ against the data.
Trade sizes come from a deterministic sample of \NTradeSizeSample{} trades from the kept markets.

\subsection{Reproducibility}\label{sec:methods-repro}
% src: S-PAPER-5, S-PAPER-6
All data queries ran on the UTRGV Cradle cluster (SLURM, CPU only).
We used Python and DuckDB \citep{raasveldt2019duckdb}.
The first-stage run is \RunId{} (job \SlurmJobId), with code version \CodeVersion.
The hash \CodeSha{} is that of the Q1 script.
The hash of the whole script set is \CodeSetSha.
The per-script hashes are in \texttt{\CodeShaFile}.
The robustness run is \RobustRun, with code version \RobustCodeVersion{} (hash \RobustCodeSha).
An earlier review found three sources of small differences between reruns.
They were tie order inside one block, a random sample, and row order.
Code version \CodeVersion{} fixes these.
The pre-registered criteria did not change, and no verdict flipped.
The review-fix run is \RThreeRun, with code version \RThreeCodeVersion{} and the pre-registration in commit \RThreePreregCommit.

\section{Results}\label{sec:results}
\subsection{Q1: is the crowd right?}\label{sec:res-q1}
% src: S-PAPER-3, S-PAPER-5, S-PAPER-10
The pre-registered rule gives this verdict: \VerdictQOne{}
The chi-square statistic is \QOneChiSq{} with \QOneDf{} degrees of freedom.
The p-value is $p=\QOneP$.
This p-value assumes independent markets, so it may be too small.
The effect is small and runs one way.
In \GapNegBinsMain{} of \GapBinsTotal{} bins the observed outcome-1 rate is below the mean price.
So outcome 1 is priced too high.
Take the \OutcomeOneYesN{} markets where outcome 1 is labeled Yes.
In \OutcomeOneYesBinsNegative{} of 9 bins (0 to 8), their rate is below the mean price.
The gaps are \OutcomeOneYesGapMin{} to \OutcomeOneYesGapMax{} percentage points.
This is descriptive, with no new test.
Over all \NMarketsKept{} markets the net gap is \GapNetDiffMain.
That is \GapNetYesShort{} fewer outcome-1 wins than the prices imply.
The largest per-bin gap is \LargestGapAbs{} (bin \LargestGapBin, counting from 0).
The largest chi-square contributor is bin \ChiTopBin{}, with \ChiTopShare{} of the chi-square.
We set no size line before the results, so small is our reading.
\Cref{fig:calib} shows the observed outcome-1 rate against the mean price in each bin.
\Cref{tab:bins} gives the numbers.

\begin{table}[H]
\centering
\caption{Calibration by price bin, with the gap as observed rate minus mean price in percent points. Why: the gap and its interval show where prices are off. Source: run \RunId.}
\label{tab:bins}
\small
\setlength{\tabcolsep}{4pt}
\begin{tabular}{lrrrcrr}
\toprule
Bin & $n$ & Mean price & Observed & \CiLevel{} CI & Gap & BH $p$ \\
\midrule
\binrow{0--10\%}{ZeroTen}\\ \midrule
\binrow{10--20\%}{TenTwenty}\\ \midrule
\binrow{20--30\%}{TwentyThirty}\\ \midrule
\binrow{30--40\%}{ThirtyForty}\\ \midrule
\binrow{40--50\%}{FortyFifty}\\ \midrule
\binrow{50--60\%}{FiftySixty}\\ \midrule
\binrow{60--70\%}{SixtySeventy}\\ \midrule
\binrow{70--80\%}{SeventyEighty}\\ \midrule
\binrow{80--90\%}{EightyNinety}\\ \midrule
\binrow{90--100\%}{NinetyHundred}\\
\bottomrule
\end{tabular}
\end{table}

The observed rate is below the mean price in every bin except the top one.
Without the minimum-trade filter the count is \GapNegBinsNoMin{} of \GapBinsTotal.
In equal-frequency bins it is \GapNegBinsEqFreq{} of \GapBinsTotal.
After BH correction, \QOneBinsBh{} of the \NBins{} bins are off ($p<\Alpha$).
After Bonferroni correction, \QOneBinsBonferroni{} bins are off.
The \CiLevel{} z interval leaves the mean price outside in \ZIntervalBinsOutsideN{} bins: \ZIntervalBinsOutsideList.
BH marks the bins \ZIntervalBinsBhOffList{} as off, which are the same bins.
For example, in the bin with mean price \BinSeventyEightyPrice{}, the observed outcome-1 rate is \BinSeventyEightyObs.
The gap is \BinSeventyEightyDiff{} (BH $p=\BinSeventyEightyBhP$).

The two cheapest bins are the longshot check.
In the cheapest bin the gap is \BinZeroTenDiff{} (BH $p=\BinZeroTenBhP$).
In the next bin the gap is \BinTenTwentyDiff{} (BH $p=\BinTenTwentyBhP$).
With the minimum-trade filter, cheap bets pay off less often than their price says.
The direction matches the horse-racing result \citep{snowberg2010favorite}.
Without the filter, the cheapest bin flips above the price.
Its gap is \NoMinBinZeroDiff{} (BH $p=\NoMinBinZeroPBh$).
The next bin has a gap of \NoMinBinOneDiff{} (BH $p=\NoMinBinOnePBh$).
Neither bin is off after BH.
So the longshot statement holds only with the filter.
Our gaps are per market and the dataset paper's are per trade.
So we compare the direction only, not the size.


\begin{figure}[H]
\centering
\figorbox{figures/fig4_calibration_gap}{\linewidth}
\caption{Observed outcome-1 rate against mean price per bin (top), and the per-bin gap in percentage points with \CiLevel{} z intervals (bottom). Why: the lower panel shows the gaps that the diagonal hides. Source: run \RunId{} and the local figure folder runs/local\_fig4gap\_v9.}
\label{fig:calib}
\end{figure}

\begin{figure}[H]
\centering
\figorbox{figures/fig5_ecdf_yes_no}{\linewidth}
\caption{ECDFs of \pOne{} for markets where outcome 1 won and markets where outcome 0 won. Why: the farther apart the ECDFs, the better the price separates the outcomes. Source: run \RunId.}
\label{fig:ecdf}
\end{figure}

Next come the score measures.
The mean squared residual of the price is \QOneBrier, the log loss is \QOneLogLoss, and the mean entropy is \QOneEntropy{} bits.
The outcome-1 rate among kept markets is \QOneYesRate, and the mean price is \PriceOneDayMean.
At the \FavoriteLine{} line, the classifier accuracy is \QOneAccuracy.
Its precision is \QOnePrecision{} and its recall is \QOneRecall.
The KS statistic between the prices of those two groups is $D=\QOneKsD$ (\Cref{fig:ecdf}).
Without the minimum-trade filter there are \QOneNoMinN{} markets.
The mean squared residual is then \QOneNoMinBrier{}, with $\chi^2=\QOneNoMinChiSq$ and $p=\QOneNoMinP$.
The chi-square test also rejects without the filter.
With \QOneEqFreqBins{} equal-frequency bins, BH marks \QOneEqFreqBinsBh{} bins as off.
The z interval leaves the price outside \ZIntervalEqFreqOutsideN{} bins, and BH marks \QOneEqFreqBinsBh.
They differ in the lowest bin, which has very few outcome-1 wins.
The z interval has no correction for the \QOneEqFreqBins{} bins.
Its spread comes from the observed rate, which is unreliable near zero.

\subsection{Q2: who is the crowd? Do a few wallets do most of the trading?}\label{sec:res-q2}
% src: S-PAPER-5, S-PAPER-10
Yes, by USDC volume, the top wallets do most of the trading.
This holds in the kept markets only, and far beyond the 80/20 reference.
The top \TopTwentyLabel{} of wallets hold \MakerTopTwentyShare{} of maker volume and \TakerTopTwentyShare{} of taker volume.
The 80/20 rule says \TopTwentyReference.
We describe these numbers without a pre-registered verdict.
Each row is one wallet and role.
Volume sums every trade in the kept markets, including trades after $T-\DaysShort$ day.
Trade-count shares were not measured.
\Cref{tab:conc} gives the top shares by role.
\Cref{fig:ccdf} shows the CCDF.
\Cref{fig:wd-maker,fig:wd-taker} show the wallet and volume distributions.
\Cref{fig:lorenz} shows the Lorenz curves.
The Lorenz curves are outside the lectures.

\begin{table}[H]
\centering
\caption{Volume concentration by role, as the share of role volume held by the top wallets. Why: the top shares show how few wallets hold most of the volume. Source: run \RunId.}
\label{tab:conc}
\small
\begin{tabular}{lrrrrrrrr}
\toprule
Role & Wallets & Mean (USDC) & Median (USDC) & Top \TopOneLabel{} & Top \TopFiveLabel{} & Top \TopTenLabel{} & Top \TopTwentyLabel \\
\midrule
Maker & \MakerNWallets & \MakerMeanVolume & \MakerMedianVolume & \MakerTopOneShare & \MakerTopFiveShare & \MakerTopTenShare & \MakerTopTwentyShare \\
\midrule
Taker & \TakerNWallets & \TakerMeanVolume & \TakerMedianVolume & \TakerTopOneShare & \TakerTopFiveShare & \TakerTopTenShare & \TakerTopTwentyShare \\
\bottomrule
\end{tabular}
\end{table}

\begin{figure}[H]
\centering
\figorbox{figures/fig6_ccdf_top_sampled}{\linewidth}
\caption{CCDF of wallet volume on log-log axes, one line per role. Why: the tail bends downward, so it is not a straight line. Source: run \CcdfRun.}
\label{fig:ccdf}
\end{figure}

\begin{figure}[H]
\centering
\begin{subfigure}[t]{0.49\linewidth}\centering
\figorbox{figures/wallet_dist_maker_count}{\linewidth}
\caption{Share of maker wallets per decade of volume.}
\label{fig:wd-maker-count}
\end{subfigure}\hfill
\begin{subfigure}[t]{0.49\linewidth}\centering
\figorbox{figures/wallet_dist_maker_volume}{\linewidth}
\caption{Share of maker volume per decade of wallet volume.}
\label{fig:wd-maker-volume}
\end{subfigure}
\caption{Wallet and volume distributions of makers, split at the top \TopTwentyLabel{} of wallets. Why: the area is the share of wallets in (a) and the share of volume in (b). Source: run \WalletDistRun.}
\label{fig:wd-maker}
\end{figure}

\begin{figure}[H]
\centering
\begin{subfigure}[t]{0.49\linewidth}\centering
\figorbox{figures/wallet_dist_taker_count}{\linewidth}
\caption{Share of taker wallets per decade of volume.}
\label{fig:wd-taker-count}
\end{subfigure}\hfill
\begin{subfigure}[t]{0.49\linewidth}\centering
\figorbox{figures/wallet_dist_taker_volume}{\linewidth}
\caption{Share of taker volume per decade of wallet volume.}
\label{fig:wd-taker-volume}
\end{subfigure}
\caption{Wallet and volume distributions of takers, split at the top \TopTwentyLabel{} of wallets. Why: the area is the share of wallets in (a) and the share of volume in (b). Source: run \WalletDistRun.}
\label{fig:wd-taker}
\end{figure}

\begin{figure}[H]
\centering
\begin{subfigure}[t]{0.49\linewidth}\centering
\figorbox{figures/wallet_lorenz_maker}{\linewidth}
\caption{Makers.}
\label{fig:lorenz-maker}
\end{subfigure}\hfill
\begin{subfigure}[t]{0.49\linewidth}\centering
\figorbox{figures/wallet_lorenz_taker}{\linewidth}
\caption{Takers.}
\label{fig:lorenz-taker}
\end{subfigure}
\caption{Lorenz curves of volume, with the 80/20 point. Why: the height is the volume share of the lowest wallets.  Source: run \WalletDistRun.}
\label{fig:lorenz}
\end{figure}

The mean volume is far above the median volume in both roles.
As a pipeline check, the dataset paper reports that the top \TopOneLabel{} of makers hold \PaperTopShare{} of maker volume \citep{qin2026polymarketv1}.
Ours is \MakerTopOneShare.
The populations differ.
The paper covers the full data, and our wallets come from the kept markets.
We use the paper's value as a reference only.

The top \TopTenLabel{} tail of wallet volume has a fitted power-law exponent $\alpha=\WalletVolumePowerAlpha$.
The fit uses \WalletVolumePowerNTail{} wallet-role rows, with both roles pooled.
The KS test rejects an exact power law ($D=\WalletVolumePowerKsD$, $p=\WalletVolumePowerKsP$).
The fit uses the same data, so the KS $p$ is too large, not too small.
So the rejection stands.
The CCDF tail in \Cref{fig:ccdf} bends downward at the largest wallets.
So it is not a straight line on log-log axes.
We do not say wallet volume follows a power law.
We set no line for how close it is.

The Welch test on log volume gives $t=\RoleWelchT$ with $p\,\RoleWelchP$, a float-underflow bound.
The effect size is the difference of mean log10 volume, makers minus takers: \RoleLogTenDiff{} (\RoleLogTenMaker{} against \RoleLogTenTaker).
The KS statistic is $D=\RoleKsD$.
Trade count and volume have a Spearman correlation of $\rho=\SpearmanTradesVolumeRho$.

In the side analysis, the fixed elbow-curve rule picks $K=\KmeansK$.
The sample has \KmeansNSample{} wallet-role rows, and a wallet in both roles gives two rows.
The \KmeansK{} clusters have \ClusterZeroN, \ClusterOneN, \ClusterTwoN, and \ClusterThreeN{} rows.
Their median volumes run from \ClusterTwoMedianVolume{} to \ClusterOneMedianVolume{} USDC.
An ANOVA of mean trade price across clusters gives $\mathrm{SSB}/(\mathrm{SSB}+\mathrm{SSW})=\KmeansPriceSsbRatio$.
Its permutation $p$ is at the floor, $1/(\KmeansNPerm+1)$.
The cluster by category test gives $\chi^2=\KmeansCategoryChi$ with \KmeansCategoryDf{} degrees of freedom.
Its $p\,\KmeansCategoryP$ is a float-underflow bound, on \KmeansNSample{} rows.
The largest cell gap is in cluster \ClusterCategoryGapCluster, category \ClusterCategoryGapCategory.
Its observed share is \ClusterCategoryGapObs{} of the sample, against \ClusterCategoryGapExp{} expected.
The gap is \ClusterCategoryGapDiff.

\subsection{Q3: what does the data look like?}\label{sec:res-q3}
% src: S-PAPER-5
We describe the distributions without a verdict.
\Cref{fig:hist} shows the price histogram and the trade-size histogram.

\begin{figure}[H]
\centering
\begin{subfigure}[t]{0.49\linewidth}\centering
\figorbox{figures/hist_price_p1d}{\linewidth}
\caption{Price \pOne{} over the kept markets.}
\label{fig:hist-price}
\end{subfigure}\hfill
\begin{subfigure}[t]{0.49\linewidth}\centering
\figorbox{figures/hist_trade_size}{\linewidth}
\caption{Trade size on a log axis. The largest sampled trade falls past the last bin edge by float rounding, so it is not plotted.}
\label{fig:hist-size}
\end{subfigure}
\caption{Histograms of the price and of trade size. Why: the price shape shows where the crowd is confident, and the trade-size shape shows how uneven bets are. Source: run \QqRun.}
\label{fig:hist}
\end{figure}

The price \pOne{} has a mean of \PriceOneDayMean{} and a median of \PriceOneDayMedian.
Prices pile up near both ends of the scale.
The trade size has a mean of \TradeSizeMean{} USDC and a median of \TradeSizeMedian{} USDC.
So trade sizes are strongly right skewed.
\Cref{fig:qq-size,fig:qq-wallet} show the QQ plots.
Points on the straight line mean the values are close to normal.

\begin{figure}[H]
\centering
\figorbox{figures/qq_trade_size_raw_log}{\linewidth}
\caption{QQ plot of trade size against the normal, raw (left) and after a log transform (right). Why: the raw points bend away from the line and the log points follow it, so the log values are closer to normal. Source: run \QqRun.}
\label{fig:qq-size}
\end{figure}

\begin{figure}[H]
\centering
\figorbox{figures/qq_wallet_volume_raw_log}{\linewidth}
\caption{QQ plot of wallet volume against the normal, raw (left) and after a log transform (right). Why: the comparison shows whether the log values are closer to normal than the raw values. Source: run \QqRun.}
\label{fig:qq-wallet}
\end{figure}

\Cref{tab:fits} gives the fit statistics.
The KS statistic $D$ is the effect size.
KS p-values are too large here because the parameters are fitted on the same data.
So a rejection stands, and a pass would not.
The normal fit is poor for both variables.
The power-law fit covers the tail only, so its $D$ is not directly comparable.
We do not rank the three fits: the normal and exponential fits use all data, and the power law uses the tail.
We did not compare a lognormal.
The exponential fit has $D=\TradeSizeExponentialKsD$ for trade size and $D=\WalletVolumeExponentialKsD$ for wallet volume.

\begin{table}[H]
\centering
\caption{Distribution fits, with the power law fitted on the tail from $x_{\min}$. Why: a small $D$ means the fitted curve is close to the data. Source: run \RunId.}
\label{tab:fits}
\small
\begin{tabular}{lrrrr}
\toprule
Variable & Power-law $\alpha$ & $x_{\min}$ (USDC) & KS $D$, power law & KS $D$, normal \\
\midrule
Trade size & \TradeSizePowerAlpha & \TradeSizePowerXmin & \TradeSizePowerKsD & \TradeSizeNormalKsD \\
\midrule
Wallet volume & \WalletVolumePowerAlpha & \WalletVolumePowerXmin & \WalletVolumePowerKsD & \WalletVolumeNormalKsD \\
\bottomrule
\end{tabular}
\end{table}

\subsection{Robustness}\label{sec:res-robust}
% src: S-PAPER-9, S-PAPER-8
Two checks look at limits of the Q1 setup.
Both are descriptive.
They add no primary test and no verdict.
The pre-registered Q1 verdict stays as it is.
The sensitivity row below sits next to the original Q1 row and does not replace it.
Both checks use the \QOneN{} kept markets.

% src: S-PAPER-9
Check R1 asks whether the Q1 result holds when markets without a clean payout are removed.
The resolution file gives each market a payout vector.
A market is clean when one outcome pays one and the other pays zero.
Of the kept markets, \RobROneCleanN{} are clean (\RobROneCleanShare).
\RobROneSplitN{} have a split payout (\RobROneSplitShare).
\RobROneOtherN{} are of another kind.
\RobROneMissingN{} is missing from the file.
Our label marks outcome 1 as the winner on \RobROneSplitYesOne{} of the split markets and outcome 0 on \RobROneSplitYesZero.
Our winner label agrees with the resolution file on \RobROneAgreement{} of the \RobROneNClean{} clean markets.
\RobROneNMismatch{} clean markets disagree.
\Cref{fig:r1} shows the payout classes and the calibration points with and without the unclear markets.
\Cref{tab:r1} puts the sensitivity row next to the original row.
We remove split, other, and mismatching markets, which is \RobROneNExcluded{} markets.
The Q1 rule result is the same in both rows: \RobROneRuleOriginal{}
The chi-square statistic moves from \QOneChiSq{} to \RobROneSensChiSq.

\begin{figure}[H]
\centering
\figorbox{figures/fig_r1_split_payout}{\linewidth}
\caption{Kept markets by payout class (log scale), and the Q1 calibration points with and without the unclear markets. Why: points that overlap mean the unclear markets do not move the result. Source: run \RobustRun.}
\label{fig:r1}
\end{figure}

\begin{table}[H]
\centering
\caption{Q1 original row and R1 sensitivity row. Why: close rows mean the unclear markets do not drive the Q1 result. Source: runs \RunId{} and \RobustRun.}
\label{tab:r1}
\small
\begin{tabular}{lrrrrr}
\toprule
Row & Markets & $\chi^2$ & $p$ & Mean sq. residual & Bins off (BH) \\
\midrule
Q1 original & \QOneN & \QOneChiSq & \QOneP & \QOneBrier & \QOneBinsBh \\
\midrule
Sensitivity & \RobROneSensN & \RobROneSensChiSq & \RobROneSensP & \RobROneSensBrier & \RobROneSensBinsBh \\
\bottomrule
\end{tabular}
\end{table}

% src: S-PAPER-9
Check R2 asks how far the price moves from $T-\DaysShort$ day to the last trade before $T$.
Call the mean price at the last timestamp before $T$ \pT.
The median absolute change from \pOne{} to \pT{} is \RobRTwoMedianAbsDelta.
The change is at least \DeltaA{} in \RobRTwoShareDeltaAtLeastTen{} of markets.
It is at least \DeltaB{} in \RobRTwoShareDeltaAtLeastTwentyFive{} and at least \DeltaC{} in \RobRTwoShareDeltaAtLeastFifty.
The price crosses \FavoriteLine{} in \RobRTwoShareCrossing{} of markets.
This compares the two end points and not the path.
\RobRTwoShareNoTrade{} of markets have no trade in the last day, and for them \pT{} equals \pOne.
\Cref{fig:r2} shows the distribution of the absolute change.

\begin{figure}[H]
\centering
\figorbox{figures/fig_r2_price_path}{\linewidth}
\caption{ECDF of the absolute price change from \pOne{} to \pT. Why: a curve near one at small values means \pOne{} is close to the last price. Source: run \RobustRun.}
\label{fig:r2}
\end{figure}

The price \pT{} is the price just before the close.
So its mean squared residual is not a one-day-ahead forecast.
The mean squared residual is \RobRTwoAllBrierOneDay{} at \pOne{} and \RobRTwoAllBrierAtClose{} at \pT.
The paired difference (\pT{} minus \pOne) is \RobRTwoAllMeanDiff{}
(\CiLevel{} CI [\RobRTwoAllCiLow, \RobRTwoAllCiHigh], $p\,\RobRTwoAllP$).
A negative difference means \pT{} is closer to the outcome than \pOne.
On the \RobRTwoTradedN{} markets with a trade in the last day, the difference is \RobRTwoTradedMeanDiff.
The \CiLevel{} CI is [\RobRTwoTradedCiLow, \RobRTwoTradedCiHigh].

\subsection{Checks requested in review}\label{sec:res-review}
% src: S-PAPER-12, S-PAPER-8
Reviewers asked for further checks.
We ran them in one run, \RThreeRun.
We fixed the rules before the run in \texttt{analysis/prereg\_review\_fixes.md}.
All checks are descriptive and add no verdict.
Each new row sits next to the original Q1 row and does not replace it.
The pre-registered Q1 verdict and its sentence stay as they are.
Every check uses the \QOneN{} kept markets, except where stated.

% src: S-PAPER-12
Check M1 asks whether the price definition drives the net gap.
Under the volume-weighted mean price (VWAP), the net gap is \RThreeMOneVwapNetGap{} (\CiLevel{} CI [\RThreeMOneVwapNetGapCiLow, \RThreeMOneVwapNetGapCiHigh]).
Under the trade-based midpoint, it is \RThreeMOneMidNetGap{} (\CiLevel{} CI [\RThreeMOneMidNetGapCiLow, \RThreeMOneMidNetGapCiHigh]).
The original row has \GapNetDiffMain.
Both intervals lie below 0.
So the pre-registered reading is that the price rule does not create the gap.
The midpoint is a trade-based proxy and not a quote midpoint.
It exists for the \RThreeMOneMidN{} markets with both trade sides (\RThreeMOneShareBoth).
The mean half spread of this proxy is $h=\RThreeMOneMeanHalfSpread$.
Is the size of the main net gap within $h$? \RThreeMOneGapWithinHalfSpread.
\Cref{fig:r3-m1} shows the three net gaps.

\begin{figure}[H]
\centering
\figorbox{figures/fig_r3_m1_price_definition}{0.8\linewidth}
\caption{Net gap under each price definition, with \CiLevel{} intervals. Why: an interval below 0 under VWAP and the midpoint means the gap is not a price-rule artifact. Source: runs \RThreeRun{} and \RThreeFigRun.}
\label{fig:r3-m1}
\end{figure}

% src: S-PAPER-12, S-PAPER-8
Check M2 splits the markets by label.
In the markets labelled Yes/No ($n=\RThreeMTwoYesNoN$), Yes is priced above its observed rate.
The net gap is \RThreeMTwoYesNoNetGap{} (\CiLevel{} CI [\RThreeMTwoYesNoNetGapCiLow, \RThreeMTwoYesNoNetGapCiHigh]).
Of \NBins{} bins, \RThreeMTwoYesNoBinsBh{} are off after BH.
The other markets show no rejection ($n=\RThreeMTwoOtherN$, $\chi^2=\RThreeMTwoOtherChiSq$, $p=\RThreeMTwoOtherP$).
Elsewhere we keep the term outcome 1.
\Cref{fig:r3-m2} shows the markets labelled Yes/No.
The figure uses Wilson intervals, which are outside the lectures.

\begin{figure}[H]
\centering
\figorbox{figures/fig_r3_m2_yes_reliability}{0.55\linewidth}
\caption{Observed Yes rate against mean Yes price in the markets labelled Yes/No. Why: points below the diagonal mean Yes is priced above its observed rate. Source: runs \RThreeRun{} and \RThreeFigRun.}
\label{fig:r3-m2}
\end{figure}

% src: S-PAPER-12
Check M3 asks how much dependence between markets matters.
The data hold no event key, so we use two proxy keys.
The reference row k0 treats each market alone, with no cluster, and is not a key.
The first key is \texttt{\RThreeMThreeReadProxy}.
It takes the first three tokens of the market slug.
It has \RThreeMThreeKOneNClusters{} clusters.
The second key is category by close hour, with \RThreeMThreeKTwoNClusters{} clusters.
We read the first key because its net gap interval is wider.
Under it, the \CiLevel{} interval is [\RThreeMThreeKOneNetGapCiLow, \RThreeMThreeKOneNetGapCiHigh], and it still lies below 0.
Under the second key it is [\RThreeMThreeKTwoNetGapCiLow, \RThreeMThreeKTwoNetGapCiHigh].
For the first key the cluster chi-square is \RThreeMThreeKOneChiSq{} ($p=\RThreeMThreeKOneP$).
For the second key it is \RThreeMThreeKTwoChiSq{} ($p=\RThreeMThreeKTwoP$).
The cluster bootstrap chi-square is outside the lectures.
The proxies are not true event keys.
\Cref{fig:r3-m3} compares the intervals.

\begin{figure}[H]
\centering
\figorbox{figures/fig_r3_m3_cluster_bootstrap}{0.8\linewidth}
\caption{Net gap with a cluster bootstrap interval under each proxy key, next to the one-market-per-cluster row. Why: a wider interval means markets in a cluster move together. Source: runs \RThreeRun{} and \RThreeFigRun.}
\label{fig:r3-m3}
\end{figure}

% src: S-PAPER-12
Check M4 compares the kept markets with the dropped ones.
The kept markets have a median volume before $T-\DaysShort$ day of \RThreeMFourKeptMedianVolumeBeforeDay{} USDC.
The pooled eligible dropped markets have \RThreeMFourDroppedMedianVolumeBeforeDay{} USDC.
The median duration is \RThreeMFourKeptMedianDuration{} against \RThreeMFourDroppedMedianDuration{} days.
The median trade count is \RThreeMFourKeptMedianTrades{} against \RThreeMFourDroppedMedianTrades.
The stale-price markets have a net gap of \RThreeMFourStaleNetGap{} (\CiLevel{} CI [\RThreeMFourStaleNetGapCiLow, \RThreeMFourStaleNetGapCiHigh]).
\Cref{fig:r3-m4} shows the share of markets in each group.

\begin{figure}[H]
\centering
\figorbox{figures/fig_r3_m4_kept_vs_dropped}{0.8\linewidth}
\caption{Share of in-scope markets in each kept and dropped group. Why: it shows how much of the market set each drop rule removes. Source: runs \RThreeRun{} and \RThreeFigRun.}
\label{fig:r3-m4}
\end{figure}

% src: S-PAPER-12
Check M5 asks how large the miscalibration is in usable terms.
We fit a logistic calibration of the outcome on the price, $\mathrm{logit}\,P(y=1)=a+b\,\mathrm{logit}(p)$.
The intercept is $a=\RThreeMFiveA$ (Wald \CiLevel{} CI [\RThreeMFiveAWaldLow, \RThreeMFiveAWaldHigh]).
The slope is $b=\RThreeMFiveB$ (Wald \CiLevel{} CI [\RThreeMFiveBWaldLow, \RThreeMFiveBWaldHigh]).
The practical line is the mean half spread, \RThreeMOneMeanHalfSpread.
The largest implied gap on the price grid is \RThreeMFiveMaxAbsDelta{} in size.
Of \RThreeMFiveGridPrices{} grid prices, \RThreeMFivePracticalPrices{} pass the practical line.
The Brier skill score against climatology is \RThreeMFiveSkill{} (Brier \RThreeMFiveBrier{} against \RThreeMFiveClimatologyBrier).
The climatology is in sample.
The majority-class accuracy is \RThreeMFiveMajorityAccuracy.
The accuracy at the \FavoriteLine{} line is \RThreeMFiveAccuracyHalf.
Logistic calibration and the skill score are outside the lectures.
\Cref{fig:r3-m5} shows the implied gap.

\begin{figure}[H]
\centering
\figorbox{figures/fig_r3_m5_implied_gap}{0.8\linewidth}
\caption{Implied gap between the fitted probability and the price, against the practical line. Why: a band that excludes 0 at a price shows the gap there is not noise. Source: runs \RThreeRun{} and \RThreeFigRun.}
\label{fig:r3-m5}
\end{figure}

% src: S-PAPER-12
The duplicate-row check removes exact duplicate trade rows.
They are \RThreeMinorFourDupShare{} of the trade rows of the kept markets.
Removing them drops \RThreeMinorFourRemoved{} markets from the kept set.
Then $\chi^2=\RThreeMinorFourChiSq$ against \QOneChiSq{} in the original row.
The net gap is \RThreeMinorFourNetGap{} against \GapNetDiffMain.
The top \TopOneLabel{} maker share changes by \RThreeMinorFourMakerTopOneDeltaPp{} percentage points.
The top \TopTwentyLabel{} maker share changes by \RThreeMinorFourMakerTopTwentyDeltaPp.
For takers the changes are \RThreeMinorFourTakerTopOneDeltaPp{} and \RThreeMinorFourTakerTopTwentyDeltaPp.

% src: S-PAPER-12
Two peer questions ask whether the gap depends on volume and on the close period.
We split the kept markets into volume terciles, by volume before $T-\DaysShort$ day.
The cut points are \RThreePeerACutLowMid{} and \RThreePeerACutMidHigh{} USDC.
The net gap is \RThreePeerALowNetGap{} in the low tercile (CI [\RThreePeerALowNetGapCiLow, \RThreePeerALowNetGapCiHigh], BH $p=\RThreePeerALowBhP$).
It is \RThreePeerAMidNetGap{} in the middle tercile (CI [\RThreePeerAMidNetGapCiLow, \RThreePeerAMidNetGapCiHigh], BH $p=\RThreePeerAMidBhP$).
It is \RThreePeerAHighNetGap{} in the high tercile (CI [\RThreePeerAHighNetGapCiLow, \RThreePeerAHighNetGapCiHigh], BH $p=\RThreePeerAHighBhP$).
The high minus low difference is \RThreePeerADiff{} (CI [\RThreePeerADiffCiLow, \RThreePeerADiffCiHigh], $p=\RThreePeerADiffP$).
\Cref{fig:r3-peer-a} shows the terciles.

\begin{figure}[H]
\centering
\figorbox{figures/fig_r3_peer_a_volume}{0.8\linewidth}
\caption{Net gap per volume tercile, with \CiLevel{} intervals. Why: the pre-registered difference test compares the high and low terciles, with high minus low \RThreePeerADiff{} (CI [\RThreePeerADiffCiLow, \RThreePeerADiffCiHigh], $p=\RThreePeerADiffP$). Source: runs \RThreeRun{} and \RThreeFigRun.}
\label{fig:r3-peer-a}
\end{figure}

The close periods are \RThreePeerBPeriodOneLabel, \RThreePeerBPeriodTwoLabel, and \RThreePeerBPeriodThreeLabel.
The net gap is \RThreePeerBPeriodOneNetGap{} in \RThreePeerBPeriodOneLabel{} (CI [\RThreePeerBPeriodOneNetGapCiLow, \RThreePeerBPeriodOneNetGapCiHigh], BH $p=\RThreePeerBPeriodOneBhP$, $n=\RThreePeerBPeriodOneN$).
It is \RThreePeerBPeriodTwoNetGap{} in \RThreePeerBPeriodTwoLabel{} (CI [\RThreePeerBPeriodTwoNetGapCiLow, \RThreePeerBPeriodTwoNetGapCiHigh], BH $p=\RThreePeerBPeriodTwoBhP$, $n=\RThreePeerBPeriodTwoN$).
In \RThreePeerBPeriodThreeLabel{} the chi-square test does not reject (BH $p=\RThreePeerBPeriodThreeBhP$, $n=\RThreePeerBPeriodThreeN$).
In the same period the net gap is \RThreePeerBPeriodThreeNetGap.
Its \CiLevel{} CI [\RThreePeerBPeriodThreeNetGapCiLow, \RThreePeerBPeriodThreeNetGapCiHigh] excludes 0.
\Cref{fig:r3-peer-b} shows the periods.

\begin{figure}[H]
\centering
\figorbox{figures/fig_r3_peer_b_period}{0.8\linewidth}
\caption{Net gap per close period, with \CiLevel{} intervals. Why: overlapping intervals mean no sign that the gap changes with the period. Source: runs \RThreeRun{} and \RThreeFigRun.}
\label{fig:r3-peer-b}
\end{figure}

\section{Discussion}\label{sec:discussion}
% src: S-PAPER-5, S-PAPER-8, S-PAPER-10
The answer to Q1 is no.
At this $n$, the price is not a calibrated probability as a whole.
The observed outcome-1 rate sits below the price in \GapNegBinsMain{} of \GapBinsTotal{} bins.
The effect is small and runs one way.
The net gap is \GapNetDiffMain.
The largest per-bin gap is \LargestGapAbs{} below the price, in bin \LargestGapBin{} (\Cref{tab:bins}).
We set no size line before the results.
With the minimum-trade filter, the longshot gap points the same way as the horse-racing bias.
Without the filter, the cheapest bin does not.
\todo{owner: interpretation of why prices run high, and what the gaps mean in practice for a reader of the price}

The answer to Q2 is descriptive.
Yes, by USDC volume, the top wallets do most of the trading.
This holds in the kept markets, in both roles, far beyond the 80/20 reference.
Our top \TopOneLabel{} maker share is \MakerTopOneShare, against \PaperTopShare{} in the dataset paper.
The populations differ.
We measure concentration of volume.
We do not test who sets the prices.
The k-means clusters are a side view and carry no claim.
The wallet volume includes trades after $T-\DaysShort$ day.

The answer to Q3 is descriptive too.
Prices cluster near the ends, and trade sizes are strongly right skewed.
\todo{owner: wider interpretation or policy implications, if the course report needs them}

\section{Limitations}\label{sec:limits}
% src: S-PAPER-1, S-PAPER-2, S-PAPER-3, S-PAPER-6, S-PAPER-10
Some markets have no usable winner label, and we drop them.
In our full run this is \NNoOutcome{} of \NMarketsInScope{} markets (\PctNoOutcome).
Our rule is wider than the dataset paper's.
It also drops markets whose outcome-1 label never traded.
The paper says only a small subset lacks a winner.

Markets keep trading after \texttt{close\_at}.
We avoid this by pricing one day before the scheduled close.
But the choice rests on a one-month pilot.

Sports is the largest topic in scope and in the kept markets.
It is \SportsShareInScope{} of the \NMarketsInScope{} markets in scope and \SportsShareKept{} of the kept markets.
Its survival is \CatSportsSurvival{} against \PctKept{} overall.
So Sports is slightly under-represented against its scope.
Sports loses markets mainly to no price (\CatSportsDropNoPriceN) or a stale price (\CatSportsDropStalePriceN).
Both are measured \DaysShort{} day before the close.
It loses almost none to the Up/Down rule (\CatSportsDropUpdownN).
The dataset paper notes that most Sports markets are Neg-Risk or Up/Down \citep{qin2026polymarketv1}.
Neg-Risk markets are outside the layer we use.

Our selection is narrow.
Only \NMarketsKept{} of \NMarketsInScope{} markets pass the filters.
\NNoPrice{} markets have no price \DaysShort{} day before the close (\PctNoPrice).
The minimum-trade filter is a choice, and we show the result without it.
\Cref{sec:res-review} compares kept and dropped markets (M4).
Markets on one event move together.
The binomial and chi-square tests assume independent markets, so the p-values are optimistic.
The z intervals use the same assumption, so they may be too narrow.
The role tests compare groups that share addresses and trades.
No event key exists in the data, so we cannot size this dependence exactly.
Two proxy keys address it in part (\Cref{sec:res-review}).
The proxies are not true event keys.
With this many rows, a tiny gap can give a tiny p-value.
So we read the effect sizes first.
The intervals do not escape the independence assumption.

The k-means result is exploratory.
The clusters describe the sample.
A wallet is an address, and one owner may control many addresses.
Volume concentration is not the same as the power to set prices.

We changed one pre-registered rule after seeing a full Q1 result.
This is the tie rule in \Cref{sec:methods-prereg}.
We report it as an amendment.
The amendment records that the verdict did not change.

The two robustness checks in \Cref{sec:res-robust} speak to two limits.
In R1, only \RobROneSplitN{} kept markets have a split payout and \RobROneMissingN{} is missing from the file.
Our label agrees with the resolution file on \RobROneAgreement{} of clean markets.
The Q1 rule result is the same without the unclear markets.
In R2, the price moves by at least \DeltaA{} in \RobRTwoShareDeltaAtLeastTen{} of markets.
This is the move between \pOne{} and the last trade.
So \pOne{} is not the price at the close.
Also \RobRTwoShareNoTrade{} of markets have no trade in the last day.

\section{Conclusion}\label{sec:conclusion}
% src: S-PAPER-5, S-PAPER-8, S-PAPER-10
We tested whether a Polymarket price can be read as a probability.
On \NMarketsKept{} resolved Yes/No markets, the answer is not exactly.
The prices are not calibrated as a whole.
With the minimum-trade filter, cheap bets pay off less often than priced.
That matches the horse-racing direction.
By USDC volume, the top wallets in the kept markets hold most of the volume.
Trade sizes are strongly right skewed.
We describe the last two results without a verdict.
The payout check leaves the Q1 result in place.
\todo{owner: closing sentence on the value of the study, and future work}

\begingroup\raggedright
\bibliographystyle{plainnat}
\bibliography{references}
\endgroup

\end{document}
